\chapter{Initial Idea and the Prototype}
\label{chp:4:InitIdeaProt}
In this chapter, a mass communication schema, and possible exceptional cases will be introduces. Afterwards, the initial prototype will be reviewed, and its drawback will be discussed.

\section{Mass Email Communication Concept}
\label{sec:4.1:MassEmaiCommConc}

When a researcher initiate a mass email communication, there will be many cases that cannot be predicted beforehand, but they affect the flow of the communication. Respondents may come up with clarifying questions to initial question, email addresses that are considered active might not be exist anymore, an auto responder might been set since the respondent is not available during a temporary period, and many more cases.

\subsection{A Mass Email Communication Scenario}
\label{subsec:4.1.1:Cust}

In section~\ref{sec:1:EmailDataCol}, the online learning platform scenario was used to illustrate the possible reasons for email data collection. If we consider it in here again, assuming that we are exploring our users to improve our platform. Therefore, we start sending out emails to our registered users to get their permission to ask some questions. This initial email might be similar to listing~\ref{lst:EmaiInvtQest}.
\vspace{1cm}

\begin{lstlisting}[language=, caption={[Email Invitation to Questionnaire for Online Learning Platform]Email Invitation to Questionnaire for Online Learning Platform}, label={lst:EmaiInvtQest}]
Dear John,

you have recently attended the "Cryptography" and the "Natural Language Processing" courses. Do you mind if you answer some questions regarding these courses?

The questionnaire will take only 15 minutes, and it will help us a lot to improve our platform.

Kind regards,
your online course team.
\end{lstlisting}

After the initial invitation email to get their permission for upcoming questionnaire. There are possible answers that we might expect:

\begin{compactenum}
	\item Yes, I would like to involve
	\item No, I am quite busy 
	\item I could involve, but I am busy till the end of this month
	\item An automated message from recipients to indicate his or her vacation period
\end{compactenum}
\vspace{1cm}

Figure~\ref{fig:EmailDiagram} shows a simplified version of this scenario. The people who accepted our invitation will get a second email including the questionnaire. As soon as they responds back with the answers of the questions, we will return a thank you email for the participation. For those who did not accept the invitation will get a motivation email to encourage for the next time. Next, some people returned conditional acceptance, which we collected them under a "maybe" case. Again, there might be several reason for them for instance they might involve in the the questionnaire but they are busy during that week. Hence, we could reply with an email telling an opinion for their situation like if people attend other exams, we can wish a good luck for it, and afterwards we can mention about a reminder email that will be send out after a time period where participants mentioned in their response. After the reminder email, we will expect to get a acceptance of the questionnaire email. Then conversation will continue as with the people who initially accepted to involve our research. The communication with auto respond case is not clear, therefore the communication will not go further with those people until they reply us in the near future. However, if the auto respond provides some information, we could also set a reminder for a upcoming date to send a reminder email.
\vspace{1cm}

\begin{figure}[htbp]
	\includegraphics[width=1.00\textwidth]{imgs/EmailDiagram.png}
	\caption[An Email Campaign Scenario]{An Email Campaign Scenario}
	\label{fig:EmailDiagram}
\end{figure}

Above mentioned sample scenario can be handled manually with individual efforts for a small group of participants. These manual effort will involve following points:

\begin{compactitem}
	\item Figure out the courses that participants attended for our email list, and add them to the content of the email for each person
	\item Write a generic answer for those who accepted to involve to the questionnaire, and write the questions
	\item Write a generic motivation email to support those who did not accept to involve to the questionnaire
	\item Go over for all the "maybe" cases, and record the date information if they mention a time period that they will be willing to participate, but busy at the moment.
	\item Write a generic email to those "maybe" cases mention about the reminder that we will send in the time period that they mentioned. This email can involve a couple sentences to give a feedback and share the feelings about respondent's situation that make them not to involve to the questionnaire at the moment.
	\item Go over for all the "auto-respond" cases to figure out if there is any condition that we need to ignore those people at all for the upcoming questionnaires in case they change their email addresses, or any date that can be recorded to setup a reminder email later on.
	\item Send a thank you email to those who replied with the answers of the questionnaire
\end{compactitem}
\vspace{1cm}

However, when the amount of participants is great, the possibility of doing this manual work will get harder, and the personalization that is added into emails will get less. For example, for the "maybe cases", we will not be able to give a feedback or share the feelings about respondent's situation, but just write an email to mention that we will send them a reminder later on. More over, these are simplified possible cases. There might be much more cases in real life situation, where it might require to ignore some respondents totally or writing a generic email by again ignoring the context of their email since individual efforts will not be enough to write custom emails for those cases. In such a case, we might even break up the communication forever, or results in the side effects of diffusion of responsibility theory as mentioned in section~\ref{sec:2.3:PersEmai}.
\vspace{1cm}

In the figure~\ref{fig:ChartEffortCustom}, you see that the personalization of emails is directly proportional to the effort to compose those emails. The shaded are on the graph represent the situation, when researchers use an application to create mass email communication, while the area marked with star represent the standard where where it is really hard to reach with the available applications. 
\vspace{1cm}

\begin{figure}[htbp]
	\centering
	\includegraphics[width=0.40\textwidth]{imgs/ChartEffortCustom.png}
	\caption[Effort vs Personalization on Composing Email]{Effort vs Personalization on Composing Email}
	\label{fig:ChartEffortCustom}
\end{figure}

In the next section, we will see that how can we help to researchers to reduce the efforts, and increase the personalization of the emails.

\subsection{Possible Improvements to Reduce Efforts}
\label{subsec:4.1.2:ImprReduEffo}

After a mass email communication scenario was described in the previous section, we will see the possible improvements in this section to reduce the efforts on the researchers' side, as well as increasing the personalization in the content in some degree.

\paragraph{Contact Information}
The first step would be to import contacts, and their relevant information. These could the information that we asked during the sign up form when they registered in our learning platform e.g. first and last name, email address, and gender. Therefore, we could easily export those information into a file format like Excel or CSV\footnote{CSV stands for comma separated values, and is a plain-text file to store tabular data to often exchange data between disparate applications \citep{Repici2010}. } to make it ready to use to initiate our mass email conversation. The application that was discussed in section~\ref{sec:3.3:Resul} offers several options to import data into those systems. Hence, researcher will not require to enter all those information manually into those systems.

\paragraph{Dynamic Variables}
The invitation email had user specific data, which are the name of the recipients at the salutation, and the courses that he or she attended. As it is hard to write those fields individually into all the emails, we could apply dynamic variables insertion approach that is used with both email marketing applications (MailChimp and Constant Contact) and one of the \ac{CRM} application (SugarCRM) analyzed in section~\ref{sec:3.3:Resul}. However, the limitation of SugarCRM was that it only allows us to record limited amount of contact profile fields. This could help us only to dynamic generation of the salutations in the emails not the course list of the recipients. On the other hand, Constant Contact and MailChimp have the feature to add custom fields into a contact profile, and dynamically use them in the emails' content. So, this could also remove the manual effort on the initial invitation email with an email template who has some dynamic fields to be replaced by the application while sending the emails. Those dynamic fields also add some degree of personalization into the emails.

\paragraph{Information Extraction}
When we start getting answers from respondents, we might need to extract some information from those emails. For example, this could be a date that we need to return to the recipients with a reminder as it is discussed in section~\ref{subsec:4.1.1:Cust}. Storing those information as we stored contact information might allow us easily access them again, and gives an opportunity to reuse them in the email content as a dynamic variable.

\paragraph{Reusability of Existing Emails}
After the initial invitation email, there will be many cases to reuse different versions of an email for further states of the conversations. We mentioned the term "template" as an email consisting of some dynamic variables in it. Later on, those variables will get their actual values, and it will be the actual message that we want to send. Here, we will add another meaning to the term "template" to emphasize a message's reusability. Let us consider the "maybe" case that was mentioned in section~\ref{subsec:4.1.1:Cust} to represent respondents who were busy to participate in the questionnaire at that moment, but maybe it is possible for them to participate in later (See figure~\ref{fig:EmailDiagram}). In those responds, recipients will come up with a date that is possible for them to participate, or an unclear answer about the availability of them. It is definite that we cannot ignore if they provide us a date, and send them a reminder before that date to participate in the questionnaire in a generic email. This will make them realize that the messages that we send are automated, and it will result the drawbacks that we discussed in chapter~\ref{chp:FouRelWor}. However, if we create two email templates: one with a dynamic variable in it to mention when we will send a reminder, and another one with a more generic content including a time far enough to encourage them to participate in the questionnaire. Therefore, when we encounter different "maybe" conditions, we have now two ready to use email templates to reply them.

\paragraph{Visualization of Conversation State}
At this point, we have began to have several email templates for different conditions to ease a researcher's life by allowing reusability on the replies. However, it has also started to difficult to pick a template to use since the state of conversation get advanced, and we have had many email templates for different conditions. This will result in additional effort to find the proper template to choose from. Besides, we will not need some of the templates after conversation's state get further, and use the latest ones more frequently than the older ones. Therefore, we need a view to show us what the current state of the conversation is, as well as allow us to pick the proper template by leveraging this view.
\vspace{1cm}

Figure~\ref{fig:ConvStatView} represents a view in a directory tree structure for this purpose. All the possible email messages that we wrote are collected in a hierarchical way starting with the first invitation email at the root node. After researcher started to get messages from respondents, all the following answers that we can write are at the second level of the tree right after the invitation email node. When ever we get an answer for the questionnaire, we write a email to say thank you. This is also represented in the third level of the tree, right after the questionnaire node.
\vspace{1cm}

\begin{figure}[htbp]
	\dirtree{%
	.1 Invitation Email.
	.2 Questionnaire - Graduate Version (GV).
	.3 Thank You GV.
	.2 Questionnaire.
	.3 Thank You.
	.2 Questionnaire After Reminder.
	.2 Next Time.
	.2 Next Time with Motivation.
	.2 Reminder with Date.
	.2 Reminder Generic.
	.2 \vdots.
	}
	\caption[A Sample Conversation State and Template View]{A Sample Conversation State and Template View}
	\label{fig:ConvStatView}
\end{figure}

A view like in figure~\ref{fig:ConvStatView} also help researchers to let them pick a template to reuse according to the email of the respondents. Therefore, they can easily see the previous answers that they gave, and the level of the tree provides an idea on the overall state of the conversation. When we consider existing applications in the market in section~\ref{sec:3.3:Resul}, none of them provided such a conversation view, and available templates had to be chosen out of mixed list which includes messages belonging other conversations as well.
\vspace{1cm}

Finally, after reviewing these initial information gathering to figure out how can we ease the efforts of the researchers, while keeping the personalization in the emails high, next section will show how this initial idea was applied into an application.


\section{Prototype}
\label{sec:4.2:Prot}
A software prototype provides developers a better understanding of the requirements, what could be developed with existing technology, and by allowing client interaction, it demonstrates them what is functionally feasible, and revise their imagination accordingly, which results better inputs from them and more forward looking systems \citep{Bernstein1996}.
\vspace{1cm}

In section~\ref{subsec:4.1.2:ImprReduEffo}, we initially considered some features to improve a researcher personalized mass email communication. Following points will put those findings into a summary to describe them as a feature of a prototype.

\begin{compactitem}
	\item Contact information can come from different sources, and assuming we have at least recipients' first names, last names, and email addresses. System should store these information under a contact list permanently. 
	\item Application user can use place holder for first and last names, therefore he or she can add those dynamic variables into email's content to personalize the salutation of the emails.
	\item Application user should extract information while reading responds into a key-value pair\footnote{A \ac{KVP} represents two conntected data items, where a key is a unique identifier for some data items, and the value can be again a data item or a pointer to the location of data \citep{Rouse2008}. In this context, this thesis uses \ac{KVP} in the context of extracted information as values, and a unique key to represents its class. For example, "first name" can be a key, and each recipients first names will be a value for the "first name" key.}.
	\item
\end{compactitem}



--> bi section da sorunlar neydi diye yaz bunla ilgili
--> Google Formla nasil kolay data aliniyor diye bahset

\begin{comment}

difficulty to pick a proper one , --> also no need to see old templates aga in since conversation got advanced


--> kimse boyle bir view sunmuyodu yaz, zendesk de var ama sadece bir kisi icin. Diger ozellikleri de saglayan app yok piyasa da diye ekleme yap

--> don't write a generic email, but personalized, effort var

--> there are many cases for maybe case
--> burada visulization koy ver graph i yapistir. listing olsun dashlerler tree yap
--> yukardakileri paragrah title'i kullnarka yaz

only supports hi part, not all

--> Improvements on the content


--> simdi neler olsa kolaytiririza basla

--> write something about this as well: http://blog.mailchimp.com/how-gmails-new-inbox-is-affecting-open-rates/
-->Merge tags allow you to dynamically add content to your email. 

--> Constant Contact'a limitation lari ekle \cite{ConstantContactInc.2011}, same for \cite{TheRocketScienceGroupLLC2013e}

Exploration
Have they taken any online courses before?

--> - Don�t forget Michael�s graph about golden standard
--> flow chart like in slides
--> - Asking follow-up questions are limited in other methods (see the book page 22: Beyond The Usability Lab - Conducting Large Scale Online User Experience Studies(2010))

%\begin{figure}
%%LaTeX with PSTricks extensions
%%Creator: inkscape 0.48.2
%%Please note this file requires PSTricks extensions
\psset{xunit=.5pt,yunit=.5pt,runit=.5pt}
\begin{pspicture}(967.5,701.25)
{
\newrgbcolor{curcolor}{1 1 1}
\pscustom[linestyle=none,fillstyle=solid,fillcolor=curcolor]
{
\newpath
\moveto(0,588.97606375)
\lineto(812.7,588.97606375)
\lineto(812.7,-0.07393625)
\lineto(0,-0.07393625)
\closepath
}
}
{
\newrgbcolor{curcolor}{0 0 0}
\pscustom[linewidth=2.10000003,linecolor=curcolor]
{
\newpath
\moveto(331.586724,587.92606375)
\lineto(474.124224,587.92606375)
\lineto(474.124224,548.02606375)
\lineto(331.586724,548.02606375)
\closepath
\moveto(331.586724,587.92606375)
}
}
{
\newrgbcolor{curcolor}{0 0 0}
\pscustom[linestyle=none,fillstyle=solid,fillcolor=curcolor]
{
\newpath
\moveto(343.46074219,573.84129812)
\lineto(346.03652344,573.84129812)
\lineto(346.03652344,563.88270438)
\lineto(343.46074219,563.88270438)
\closepath
\moveto(343.46074219,573.84129812)
}
}
{
\newrgbcolor{curcolor}{0 0 0}
\pscustom[linestyle=none,fillstyle=solid,fillcolor=curcolor]
{
\newpath
\moveto(355.94589844,568.42723562)
\lineto(355.94589844,563.88270438)
\lineto(353.55058594,563.88270438)
\lineto(353.55058594,567.36082937)
\curveto(353.55058594,568.00477443)(353.53417969,568.44774318)(353.50136719,568.68973563)
\curveto(353.47675729,568.93993068)(353.42753959,569.12039942)(353.35371094,569.23114188)
\curveto(353.26347709,569.38289943)(353.14042969,569.50184447)(352.97636719,569.59207938)
\curveto(352.82050729,569.67821193)(352.64003959,569.72332938)(352.43496094,569.72332938)
\curveto(351.93046849,569.72332938)(351.53671849,569.53055567)(351.25371094,569.14911063)
\curveto(350.96660209,568.76356428)(350.82714844,568.23446193)(350.82714844,567.55770437)
\lineto(350.82714844,563.88270438)
\lineto(348.44824219,563.88270438)
\lineto(348.44824219,571.34754813)
\lineto(350.82714844,571.34754813)
\lineto(350.82714844,570.24832938)
\curveto(351.18808594,570.68309447)(351.56953099,571.00711817)(351.97558594,571.21629812)
\curveto(352.37753959,571.42137677)(352.82871094,571.52801688)(353.32089844,571.52801688)
\curveto(354.18222709,571.52801688)(354.83437474,571.25731427)(355.27324219,570.72411062)
\curveto(355.72031224,570.19911063)(355.94589844,569.43211818)(355.94589844,568.42723562)
\closepath
\moveto(355.94589844,568.42723562)
}
}
{
\newrgbcolor{curcolor}{0 0 0}
\pscustom[linestyle=none,fillstyle=solid,fillcolor=curcolor]
{
\newpath
\moveto(357.23378906,571.34754813)
\lineto(359.61269531,571.34754813)
\lineto(361.46660156,566.19598563)
\lineto(363.32050781,571.34754813)
\lineto(365.71582031,571.34754813)
\lineto(362.77910156,563.88270438)
\lineto(360.15410156,563.88270438)
\closepath
\moveto(357.23378906,571.34754813)
}
}
{
\newrgbcolor{curcolor}{0 0 0}
\pscustom[linestyle=none,fillstyle=solid,fillcolor=curcolor]
{
\newpath
\moveto(367.06933594,571.34754813)
\lineto(369.44824219,571.34754813)
\lineto(369.44824219,563.88270438)
\lineto(367.06933594,563.88270438)
\closepath
\moveto(367.06933594,574.25145437)
\lineto(369.44824219,574.25145437)
\lineto(369.44824219,572.31551687)
\lineto(367.06933594,572.31551687)
\closepath
\moveto(367.06933594,574.25145437)
}
}
{
\newrgbcolor{curcolor}{0 0 0}
\pscustom[linestyle=none,fillstyle=solid,fillcolor=curcolor]
{
\newpath
\moveto(374.35371094,573.46395437)
\lineto(374.35371094,571.34754813)
\lineto(376.81464844,571.34754813)
\lineto(376.81464844,569.64129813)
\lineto(374.35371094,569.64129813)
\lineto(374.35371094,566.47489188)
\curveto(374.35371094,566.12215698)(374.41933594,565.88836817)(374.55058594,565.76942313)
\curveto(374.69003959,565.64637678)(374.96484349,565.58895437)(375.37089844,565.58895437)
\lineto(376.60136719,565.58895437)
\lineto(376.60136719,563.88270438)
\lineto(374.55058594,563.88270438)
\curveto(373.60722709,563.88270438)(372.94277344,564.07957937)(372.54902344,564.47332938)
\curveto(372.15527344,564.86707938)(371.95839844,565.53153198)(371.95839844,566.47489188)
\lineto(371.95839844,569.64129813)
\lineto(370.77714844,569.64129813)
\lineto(370.77714844,571.34754813)
\lineto(371.95839844,571.34754813)
\lineto(371.95839844,573.46395437)
\closepath
\moveto(374.35371094,573.46395437)
}
}
{
\newrgbcolor{curcolor}{0 0 0}
\pscustom[linestyle=none,fillstyle=solid,fillcolor=curcolor]
{
\newpath
\moveto(381.61347656,567.24598563)
\curveto(381.10898411,567.24598563)(380.73164036,567.15575177)(380.48144531,566.98348562)
\curveto(380.22714896,566.81942312)(380.10410156,566.57332937)(380.10410156,566.24520437)
\curveto(380.10410156,565.93758693)(380.20253906,565.69559448)(380.39941406,565.52332938)
\curveto(380.60449166,565.35926688)(380.89160156,565.27723563)(381.25253906,565.27723563)
\curveto(381.69960911,565.27723563)(382.07695286,565.43309447)(382.38457031,565.75301687)
\curveto(382.68808646,566.08114187)(382.84394531,566.48309448)(382.84394531,566.96707938)
\lineto(382.84394531,567.24598563)
\closepath
\moveto(385.25566406,568.14832938)
\lineto(385.25566406,563.88270438)
\lineto(382.84394531,563.88270438)
\lineto(382.84394531,564.98192312)
\curveto(382.52402396,564.53075177)(382.16308646,564.20262677)(381.76113281,563.99754812)
\curveto(381.36738281,563.79247)(380.88339896,563.68582938)(380.31738281,563.68582938)
\curveto(379.53808646,563.68582938)(378.91054661,563.90731375)(378.43066406,564.35848562)
\curveto(377.94667916,564.81786062)(377.70878906,565.40848563)(377.70878906,566.13036062)
\curveto(377.70878906,567.00399318)(378.00820286,567.64383692)(378.61113281,568.04989187)
\curveto(379.20996146,568.45184448)(380.15742161,568.65692313)(381.44941406,568.65692313)
\lineto(382.84394531,568.65692313)
\lineto(382.84394531,568.85379813)
\curveto(382.84394531,569.22293928)(382.69628906,569.49774317)(382.40097656,569.67411063)
\curveto(382.10566406,569.84637677)(381.63808646,569.93661062)(381.00644531,569.93661062)
\curveto(380.50195286,569.93661062)(380.02617161,569.88739188)(379.57910156,569.78895438)
\curveto(379.14023411,569.69051687)(378.73417916,569.53465697)(378.36503906,569.32957937)
\lineto(378.36503906,571.15067313)
\curveto(378.86542916,571.26961818)(379.37402396,571.36395438)(379.89082031,571.42957938)
\curveto(380.41582031,571.49520437)(380.93261666,571.52801688)(381.44941406,571.52801688)
\curveto(382.79472656,571.52801688)(383.76269531,571.25731427)(384.35332031,570.72411062)
\curveto(384.95214896,570.19911063)(385.25566406,569.33778197)(385.25566406,568.14832938)
\closepath
\moveto(385.25566406,568.14832938)
}
}
{
\newrgbcolor{curcolor}{0 0 0}
\pscustom[linestyle=none,fillstyle=solid,fillcolor=curcolor]
{
\newpath
\moveto(390.08320313,573.46395437)
\lineto(390.08320313,571.34754813)
\lineto(392.54414063,571.34754813)
\lineto(392.54414063,569.64129813)
\lineto(390.08320313,569.64129813)
\lineto(390.08320313,566.47489188)
\curveto(390.08320313,566.12215698)(390.14882813,565.88836817)(390.28007812,565.76942313)
\curveto(390.41953177,565.64637678)(390.69433567,565.58895437)(391.10039062,565.58895437)
\lineto(392.33085937,565.58895437)
\lineto(392.33085937,563.88270438)
\lineto(390.28007812,563.88270438)
\curveto(389.33671928,563.88270438)(388.67226563,564.07957937)(388.27851563,564.47332938)
\curveto(387.88476562,564.86707938)(387.68789063,565.53153198)(387.68789063,566.47489188)
\lineto(387.68789063,569.64129813)
\lineto(386.50664062,569.64129813)
\lineto(386.50664062,571.34754813)
\lineto(387.68789063,571.34754813)
\lineto(387.68789063,573.46395437)
\closepath
\moveto(390.08320313,573.46395437)
}
}
{
\newrgbcolor{curcolor}{0 0 0}
\pscustom[linestyle=none,fillstyle=solid,fillcolor=curcolor]
{
\newpath
\moveto(393.99609375,571.34754813)
\lineto(396.375,571.34754813)
\lineto(396.375,563.88270438)
\lineto(393.99609375,563.88270438)
\closepath
\moveto(393.99609375,574.25145437)
\lineto(396.375,574.25145437)
\lineto(396.375,572.31551687)
\lineto(393.99609375,572.31551687)
\closepath
\moveto(393.99609375,574.25145437)
}
}
{
\newrgbcolor{curcolor}{0 0 0}
\pscustom[linestyle=none,fillstyle=solid,fillcolor=curcolor]
{
\newpath
\moveto(402.215625,569.82176687)
\curveto(401.690625,569.82176687)(401.28457005,569.62899318)(401.0015625,569.24754812)
\curveto(400.72675755,568.87430567)(400.59140625,568.32879813)(400.59140625,567.60692312)
\curveto(400.59140625,566.89325177)(400.72675755,566.34774318)(401.0015625,565.96629812)
\curveto(401.28457005,565.58075177)(401.690625,565.39207937)(402.215625,565.39207937)
\curveto(402.740625,565.39207937)(403.1384763,565.58075177)(403.41328125,565.96629812)
\curveto(403.6839849,566.34774318)(403.8234375,566.89325177)(403.8234375,567.60692312)
\curveto(403.8234375,568.32879813)(403.6839849,568.87430567)(403.41328125,569.24754812)
\curveto(403.1384763,569.62899318)(402.740625,569.82176687)(402.215625,569.82176687)
\closepath
\moveto(402.215625,571.52801688)
\curveto(403.5035151,571.52801688)(404.5125,571.17528197)(405.234375,570.47801687)
\curveto(405.95625,569.78895438)(406.3171875,568.82918927)(406.3171875,567.60692312)
\curveto(406.3171875,566.38055567)(405.95625,565.41668927)(405.234375,564.71942313)
\curveto(404.5125,564.03036062)(403.5035151,563.68582938)(402.215625,563.68582938)
\curveto(400.9359375,563.68582938)(399.9269526,564.03036062)(399.196875,564.71942313)
\curveto(398.475,565.41668927)(398.1140625,566.38055567)(398.1140625,567.60692312)
\curveto(398.1140625,568.82918927)(398.475,569.78895438)(399.196875,570.47801687)
\curveto(399.9269526,571.17528197)(400.9359375,571.52801688)(402.215625,571.52801688)
\closepath
\moveto(402.215625,571.52801688)
}
}
{
\newrgbcolor{curcolor}{0 0 0}
\pscustom[linestyle=none,fillstyle=solid,fillcolor=curcolor]
{
\newpath
\moveto(415.54160156,568.42723562)
\lineto(415.54160156,563.88270438)
\lineto(413.14628906,563.88270438)
\lineto(413.14628906,567.36082937)
\curveto(413.14628906,568.00477443)(413.12988281,568.44774318)(413.09707031,568.68973563)
\curveto(413.07246041,568.93993068)(413.02324271,569.12039942)(412.94941406,569.23114188)
\curveto(412.85918021,569.38289943)(412.73613281,569.50184447)(412.57207031,569.59207938)
\curveto(412.41621041,569.67821193)(412.23574271,569.72332938)(412.03066406,569.72332938)
\curveto(411.52617161,569.72332938)(411.13242161,569.53055567)(410.84941406,569.14911063)
\curveto(410.56230521,568.76356428)(410.42285156,568.23446193)(410.42285156,567.55770437)
\lineto(410.42285156,563.88270438)
\lineto(408.04394531,563.88270438)
\lineto(408.04394531,571.34754813)
\lineto(410.42285156,571.34754813)
\lineto(410.42285156,570.24832938)
\curveto(410.78378906,570.68309447)(411.16523411,571.00711817)(411.57128906,571.21629812)
\curveto(411.97324271,571.42137677)(412.42441406,571.52801688)(412.91660156,571.52801688)
\curveto(413.77793021,571.52801688)(414.43007786,571.25731427)(414.86894531,570.72411062)
\curveto(415.31601536,570.19911063)(415.54160156,569.43211818)(415.54160156,568.42723562)
\closepath
\moveto(415.54160156,568.42723562)
}
}
{
\newrgbcolor{curcolor}{0 0 0}
\pscustom[linestyle=none,fillstyle=solid,fillcolor=curcolor]
{
\newpath
\moveto(422.62089844,573.84129812)
\lineto(429.54433594,573.84129812)
\lineto(429.54433594,571.88895438)
\lineto(425.19667969,571.88895438)
\lineto(425.19667969,570.03504813)
\lineto(429.29824219,570.03504813)
\lineto(429.29824219,568.09911062)
\lineto(425.19667969,568.09911062)
\lineto(425.19667969,565.81864188)
\lineto(429.69199219,565.81864188)
\lineto(429.69199219,563.88270438)
\lineto(422.62089844,563.88270438)
\closepath
\moveto(422.62089844,573.84129812)
}
}
{
\newrgbcolor{curcolor}{0 0 0}
\pscustom[linestyle=none,fillstyle=solid,fillcolor=curcolor]
{
\newpath
\moveto(438.77695312,570.10067312)
\curveto(439.07226562,570.56825178)(439.42500052,570.92508692)(439.84335937,571.16707937)
\curveto(440.25761692,571.40496948)(440.71699193,571.52801688)(441.22148437,571.52801688)
\curveto(442.07460938,571.52801688)(442.72265572,571.25731427)(443.17382812,570.72411062)
\curveto(443.62089817,570.19911063)(443.84648437,569.43211818)(443.84648437,568.42723562)
\lineto(443.84648437,563.88270438)
\lineto(441.45117188,563.88270438)
\lineto(441.45117188,567.77098562)
\lineto(441.45117188,567.95145438)
\curveto(441.45937552,568.01707937)(441.46757812,568.10321192)(441.46757812,568.21395437)
\curveto(441.46757812,568.74715698)(441.38554687,569.13270438)(441.22148437,569.36239187)
\curveto(441.06562552,569.60028198)(440.81953072,569.72332938)(440.48320312,569.72332938)
\curveto(440.02382813,569.72332938)(439.66699192,569.53465697)(439.41679687,569.16551687)
\curveto(439.17480442,568.79227442)(439.04765573,568.25086818)(439.03945313,567.54129813)
\lineto(439.03945313,563.88270438)
\lineto(436.64414062,563.88270438)
\lineto(436.64414062,567.77098562)
\curveto(436.64414062,568.59950178)(436.57031303,569.13270438)(436.43085938,569.36239187)
\curveto(436.28730443,569.60028198)(436.03710938,569.72332938)(435.67617188,569.72332938)
\curveto(435.21679688,569.72332938)(434.85996068,569.53465697)(434.60976563,569.16551687)
\curveto(434.35546822,568.79227442)(434.23242188,568.25496948)(434.23242188,567.55770437)
\lineto(434.23242188,563.88270438)
\lineto(431.83710938,563.88270438)
\lineto(431.83710938,571.34754813)
\lineto(434.23242188,571.34754813)
\lineto(434.23242188,570.24832938)
\curveto(434.52773438,570.67489187)(434.86406303,570.99071193)(435.24960937,571.19989188)
\curveto(435.63105443,571.41727442)(436.04531302,571.52801688)(436.49648438,571.52801688)
\curveto(437.02148438,571.52801688)(437.48085938,571.40086818)(437.87460937,571.15067313)
\curveto(438.26835937,570.89637678)(438.56777318,570.54774318)(438.77695312,570.10067312)
\closepath
\moveto(438.77695312,570.10067312)
}
}
{
\newrgbcolor{curcolor}{0 0 0}
\pscustom[linestyle=none,fillstyle=solid,fillcolor=curcolor]
{
\newpath
\moveto(449.4328125,567.24598563)
\curveto(448.92832005,567.24598563)(448.5509763,567.15575177)(448.30078125,566.98348562)
\curveto(448.0464849,566.81942312)(447.9234375,566.57332937)(447.9234375,566.24520437)
\curveto(447.9234375,565.93758693)(448.021875,565.69559448)(448.21875,565.52332938)
\curveto(448.4238276,565.35926688)(448.7109375,565.27723563)(449.071875,565.27723563)
\curveto(449.51894505,565.27723563)(449.8962888,565.43309447)(450.20390625,565.75301687)
\curveto(450.5074224,566.08114187)(450.66328125,566.48309448)(450.66328125,566.96707938)
\lineto(450.66328125,567.24598563)
\closepath
\moveto(453.075,568.14832938)
\lineto(453.075,563.88270438)
\lineto(450.66328125,563.88270438)
\lineto(450.66328125,564.98192312)
\curveto(450.3433599,564.53075177)(449.9824224,564.20262677)(449.58046875,563.99754812)
\curveto(449.18671875,563.79247)(448.7027349,563.68582938)(448.13671875,563.68582938)
\curveto(447.3574224,563.68582938)(446.72988255,563.90731375)(446.25,564.35848562)
\curveto(445.7660151,564.81786062)(445.528125,565.40848563)(445.528125,566.13036062)
\curveto(445.528125,567.00399318)(445.8275388,567.64383692)(446.43046875,568.04989187)
\curveto(447.0292974,568.45184448)(447.97675755,568.65692313)(449.26875,568.65692313)
\lineto(450.66328125,568.65692313)
\lineto(450.66328125,568.85379813)
\curveto(450.66328125,569.22293928)(450.515625,569.49774317)(450.2203125,569.67411063)
\curveto(449.925,569.84637677)(449.4574224,569.93661062)(448.82578125,569.93661062)
\curveto(448.3212888,569.93661062)(447.84550755,569.88739188)(447.3984375,569.78895438)
\curveto(446.95957005,569.69051687)(446.5535151,569.53465697)(446.184375,569.32957937)
\lineto(446.184375,571.15067313)
\curveto(446.6847651,571.26961818)(447.1933599,571.36395438)(447.71015625,571.42957938)
\curveto(448.23515625,571.49520437)(448.7519526,571.52801688)(449.26875,571.52801688)
\curveto(450.6140625,571.52801688)(451.58203125,571.25731427)(452.17265625,570.72411062)
\curveto(452.7714849,570.19911063)(453.075,569.33778197)(453.075,568.14832938)
\closepath
\moveto(453.075,568.14832938)
}
}
{
\newrgbcolor{curcolor}{0 0 0}
\pscustom[linestyle=none,fillstyle=solid,fillcolor=curcolor]
{
\newpath
\moveto(455.29394531,571.34754813)
\lineto(457.67285156,571.34754813)
\lineto(457.67285156,563.88270438)
\lineto(455.29394531,563.88270438)
\closepath
\moveto(455.29394531,574.25145437)
\lineto(457.67285156,574.25145437)
\lineto(457.67285156,572.31551687)
\lineto(455.29394531,572.31551687)
\closepath
\moveto(455.29394531,574.25145437)
}
}
{
\newrgbcolor{curcolor}{0 0 0}
\pscustom[linestyle=none,fillstyle=solid,fillcolor=curcolor]
{
\newpath
\moveto(459.96972656,574.25145437)
\lineto(462.34863281,574.25145437)
\lineto(462.34863281,563.88270438)
\lineto(459.96972656,563.88270438)
\closepath
\moveto(459.96972656,574.25145437)
}
}
{
\newrgbcolor{curcolor}{0 0 0}
\pscustom[linewidth=2.10000003,linecolor=curcolor]
{
\newpath
\moveto(812.179095,433.57606375)
\curveto(812.179095,414.35203375)(781.00722,398.76610675)(742.559181,398.76610675)
\curveto(704.107026,398.76610675)(672.935151,414.35203375)(672.935151,433.57606375)
\curveto(672.935151,452.80009375)(704.107026,468.38602075)(742.559181,468.38602075)
\curveto(781.00722,468.38602075)(812.179095,452.80009375)(812.179095,433.57606375)
}
}
{
\newrgbcolor{curcolor}{0 0 0}
\pscustom[linestyle=none,fillstyle=solid,fillcolor=curcolor]
{
\newpath
\moveto(731.75566406,439.69969656)
\lineto(727.75253906,439.69969656)
\lineto(727.11269531,437.87860281)
\lineto(724.53691406,437.87860281)
\lineto(728.22832031,447.83719656)
\lineto(731.27988281,447.83719656)
\lineto(734.97128906,437.87860281)
\lineto(732.39550781,437.87860281)
\closepath
\moveto(728.39238281,441.53719656)
\lineto(731.11582031,441.53719656)
\lineto(729.75410156,445.49110281)
\closepath
\moveto(728.39238281,441.53719656)
}
}
{
\newrgbcolor{curcolor}{0 0 0}
\pscustom[linestyle=none,fillstyle=solid,fillcolor=curcolor]
{
\newpath
\moveto(736.09921875,440.78250906)
\lineto(736.09921875,445.34344656)
\lineto(738.49453125,445.34344656)
\lineto(738.49453125,444.60516531)
\curveto(738.49453125,444.19911036)(738.4863276,443.69051661)(738.478125,443.07938406)
\lineto(738.478125,441.84891531)
\curveto(738.478125,441.24598536)(738.49453125,440.81532156)(738.52734375,440.55282156)
\curveto(738.56015625,440.29032156)(738.6134763,440.09754786)(738.69140625,439.97860281)
\curveto(738.78984375,439.82274291)(738.9128901,439.70379786)(739.06875,439.61766531)
\curveto(739.2328125,439.52743041)(739.41738255,439.48641531)(739.6265625,439.48641531)
\curveto(740.11875,439.48641531)(740.5042974,439.67508771)(740.79140625,440.06063406)
\curveto(741.0744138,440.44207911)(741.21796875,440.97118041)(741.21796875,441.65204031)
\lineto(741.21796875,445.34344656)
\lineto(743.596875,445.34344656)
\lineto(743.596875,437.87860281)
\lineto(741.21796875,437.87860281)
\lineto(741.21796875,438.96141531)
\curveto(740.85703125,438.52254786)(740.4714849,438.19852521)(740.06953125,437.99344656)
\curveto(739.67578125,437.78836844)(739.2328125,437.68172781)(738.740625,437.68172781)
\curveto(737.87519505,437.68172781)(737.21894505,437.94422781)(736.771875,438.46922781)
\curveto(736.3207026,439.00243041)(736.09921875,439.77352521)(736.09921875,440.78250906)
\closepath
\moveto(736.09921875,440.78250906)
}
}
{
\newrgbcolor{curcolor}{0 0 0}
\pscustom[linestyle=none,fillstyle=solid,fillcolor=curcolor]
{
\newpath
\moveto(748.51054687,447.45985281)
\lineto(748.51054687,445.34344656)
\lineto(750.97148437,445.34344656)
\lineto(750.97148437,443.63719656)
\lineto(748.51054687,443.63719656)
\lineto(748.51054687,440.47079031)
\curveto(748.51054687,440.11805541)(748.57617188,439.88426661)(748.70742188,439.76532156)
\curveto(748.84687552,439.64227521)(749.12167943,439.58485281)(749.52773438,439.58485281)
\lineto(750.75820313,439.58485281)
\lineto(750.75820313,437.87860281)
\lineto(748.70742188,437.87860281)
\curveto(747.76406302,437.87860281)(747.09960938,438.07547781)(746.70585937,438.46922781)
\curveto(746.31210937,438.86297781)(746.11523438,439.52743041)(746.11523438,440.47079031)
\lineto(746.11523438,443.63719656)
\lineto(744.93398438,443.63719656)
\lineto(744.93398438,445.34344656)
\lineto(746.11523438,445.34344656)
\lineto(746.11523438,447.45985281)
\closepath
\moveto(748.51054687,447.45985281)
}
}
{
\newrgbcolor{curcolor}{0 0 0}
\pscustom[linestyle=none,fillstyle=solid,fillcolor=curcolor]
{
\newpath
\moveto(755.9671875,443.81766531)
\curveto(755.4421875,443.81766531)(755.03613255,443.62489161)(754.753125,443.24344656)
\curveto(754.47832005,442.87020411)(754.34296875,442.32469656)(754.34296875,441.60282156)
\curveto(754.34296875,440.88915021)(754.47832005,440.34364161)(754.753125,439.96219656)
\curveto(755.03613255,439.57665021)(755.4421875,439.38797781)(755.9671875,439.38797781)
\curveto(756.4921875,439.38797781)(756.8900388,439.57665021)(757.16484375,439.96219656)
\curveto(757.4355474,440.34364161)(757.575,440.88915021)(757.575,441.60282156)
\curveto(757.575,442.32469656)(757.4355474,442.87020411)(757.16484375,443.24344656)
\curveto(756.8900388,443.62489161)(756.4921875,443.81766531)(755.9671875,443.81766531)
\closepath
\moveto(755.9671875,445.52391531)
\curveto(757.2550776,445.52391531)(758.2640625,445.17118041)(758.9859375,444.47391531)
\curveto(759.7078125,443.78485281)(760.06875,442.82508771)(760.06875,441.60282156)
\curveto(760.06875,440.37645411)(759.7078125,439.41258771)(758.9859375,438.71532156)
\curveto(758.2640625,438.02625906)(757.2550776,437.68172781)(755.9671875,437.68172781)
\curveto(754.6875,437.68172781)(753.6785151,438.02625906)(752.9484375,438.71532156)
\curveto(752.2265625,439.41258771)(751.865625,440.37645411)(751.865625,441.60282156)
\curveto(751.865625,442.82508771)(752.2265625,443.78485281)(752.9484375,444.47391531)
\curveto(753.6785151,445.17118041)(754.6875,445.52391531)(755.9671875,445.52391531)
\closepath
\moveto(755.9671875,445.52391531)
}
}
{
\newrgbcolor{curcolor}{0 0 0}
\pscustom[linestyle=none,fillstyle=solid,fillcolor=curcolor]
{
\newpath
\moveto(706.2234375,426.62391531)
\curveto(706.75664115,426.62391531)(707.1421875,426.72235281)(707.371875,426.91922781)
\curveto(707.6015625,427.11610281)(707.71640625,427.44422781)(707.71640625,427.90360281)
\curveto(707.71640625,428.35067286)(707.6015625,428.67469656)(707.371875,428.87157156)
\curveto(707.1421875,429.06844656)(706.75664115,429.16688406)(706.2234375,429.16688406)
\lineto(705.140625,429.16688406)
\lineto(705.140625,426.62391531)
\closepath
\moveto(705.140625,424.85204031)
\lineto(705.140625,421.07860281)
\lineto(702.56484375,421.07860281)
\lineto(702.56484375,431.03719656)
\lineto(706.4859375,431.03719656)
\curveto(707.7984375,431.03719656)(708.75820365,430.81161036)(709.3734375,430.36454031)
\curveto(709.98457005,429.92567286)(710.2921875,429.23250906)(710.2921875,428.28094656)
\curveto(710.2921875,427.62469656)(710.1322263,427.08329031)(709.81640625,426.65672781)
\curveto(709.49648385,426.23836791)(709.0166013,425.93485281)(708.37265625,425.73797781)
\curveto(708.7212888,425.64774291)(709.03300755,425.46317286)(709.3078125,425.18016531)
\curveto(709.59082005,424.90536036)(709.87382865,424.47879786)(710.1609375,423.90047781)
\lineto(711.55546875,421.07860281)
\lineto(708.815625,421.07860281)
\lineto(707.61796875,423.55594656)
\curveto(707.36367135,424.04813406)(707.1134763,424.38446271)(706.86328125,424.57313406)
\curveto(706.60898385,424.75770411)(706.27675755,424.85204031)(705.8625,424.85204031)
\closepath
\moveto(705.140625,424.85204031)
}
}
{
\newrgbcolor{curcolor}{0 0 0}
\pscustom[linestyle=none,fillstyle=solid,fillcolor=curcolor]
{
\newpath
\moveto(720.43535156,424.83563406)
\lineto(720.43535156,424.14657156)
\lineto(714.85722656,424.14657156)
\curveto(714.91054661,423.58875906)(715.11152396,423.16629786)(715.46425781,422.88329031)
\curveto(715.81289036,422.60848536)(716.30097656,422.47313406)(716.92441406,422.47313406)
\curveto(717.41660156,422.47313406)(717.92519531,422.54286036)(718.45019531,422.68641531)
\curveto(718.98339896,422.83817286)(719.53300781,423.06786036)(720.09082031,423.37547781)
\lineto(720.09082031,421.53797781)
\curveto(719.52070286,421.31649291)(718.95058646,421.15243094)(718.38457031,421.04579031)
\curveto(717.82675781,420.93915021)(717.26074166,420.88172781)(716.69472656,420.88172781)
\curveto(715.34941406,420.88172781)(714.29941406,421.22625906)(713.54472656,421.91532156)
\curveto(712.79824166,422.60438406)(712.42910156,423.56415021)(712.42910156,424.80282156)
\curveto(712.42910156,426.01688406)(712.79414036,426.97254786)(713.52832031,427.67391531)
\curveto(714.25839896,428.37118041)(715.27148411,428.72391531)(716.56347656,428.72391531)
\curveto(717.73242161,428.72391531)(718.66757786,428.36707911)(719.36894531,427.65750906)
\curveto(720.07851536,426.95614161)(720.43535156,426.01688406)(720.43535156,424.83563406)
\closepath
\moveto(717.99082031,425.62313406)
\curveto(717.99082031,426.08250906)(717.85136666,426.44754786)(717.58066406,426.72235281)
\curveto(717.31816406,426.99305541)(716.97363281,427.13250906)(716.54707031,427.13250906)
\curveto(716.08769531,427.13250906)(715.71035156,427.00125906)(715.41503906,426.73875906)
\curveto(715.12792916,426.48446271)(714.95566406,426.11532156)(714.89003906,425.62313406)
\closepath
\moveto(717.99082031,425.62313406)
}
}
{
\newrgbcolor{curcolor}{0 0 0}
\pscustom[linestyle=none,fillstyle=solid,fillcolor=curcolor]
{
\newpath
\moveto(728.06015625,428.31375906)
\lineto(728.06015625,426.49266531)
\curveto(727.5556638,426.71004786)(727.0634763,426.87000906)(726.58359375,426.96844656)
\curveto(726.1119138,427.07508771)(725.66484375,427.13250906)(725.23828125,427.13250906)
\curveto(724.7871099,427.13250906)(724.45078125,427.07508771)(724.22109375,426.96844656)
\curveto(723.9996099,426.85770411)(723.89296875,426.68133771)(723.89296875,426.44344656)
\curveto(723.89296875,426.25477521)(723.975,426.10711791)(724.1390625,426.00047781)
\curveto(724.3113276,425.90204031)(724.61484375,425.82411036)(725.04140625,425.77079031)
\lineto(725.46796875,425.72157156)
\curveto(726.6902349,425.56571271)(727.5105474,425.31141531)(727.92890625,424.95047781)
\curveto(728.35546875,424.59774291)(728.56875,424.03993041)(728.56875,423.27704031)
\curveto(728.56875,422.47723536)(728.2734375,421.87430541)(727.6828125,421.47235281)
\curveto(727.1003901,421.07860281)(726.22675755,420.88172781)(725.0578125,420.88172781)
\curveto(724.565625,420.88172781)(724.05703125,420.92274291)(723.53203125,420.99657156)
\curveto(723.00703125,421.07450125)(722.465625,421.19344656)(721.9078125,421.35750906)
\lineto(721.9078125,423.16219656)
\curveto(722.3753901,422.93250906)(722.8634763,422.75614161)(723.36796875,422.63719656)
\curveto(723.8683599,422.52645411)(724.3769526,422.47313406)(724.89375,422.47313406)
\curveto(725.3613276,422.47313406)(725.7140625,422.53055541)(725.94375,422.65360281)
\curveto(726.1816401,422.78485281)(726.3046875,422.98172781)(726.3046875,423.24422781)
\curveto(726.3046875,423.44930541)(726.22265625,423.60516531)(726.05859375,423.70360281)
\curveto(725.89453125,423.81024291)(725.57050755,423.90047781)(725.090625,423.96610281)
\lineto(724.6640625,424.01532156)
\curveto(723.60175755,424.14657156)(722.859375,424.39266531)(722.4328125,424.75360281)
\curveto(722.00625,425.11454031)(721.79296875,425.66004786)(721.79296875,426.39422781)
\curveto(721.79296875,427.18172781)(722.05957005,427.76415021)(722.596875,428.14969656)
\curveto(723.14238255,428.53114161)(723.9791013,428.72391531)(725.10703125,428.72391531)
\curveto(725.5417974,428.72391531)(726.0011724,428.69110281)(726.48515625,428.62547781)
\curveto(726.9650388,428.55985281)(727.4900388,428.45321271)(728.06015625,428.31375906)
\closepath
\moveto(728.06015625,428.31375906)
}
}
{
\newrgbcolor{curcolor}{0 0 0}
\pscustom[linestyle=none,fillstyle=solid,fillcolor=curcolor]
{
\newpath
\moveto(732.7359375,422.16141531)
\lineto(732.7359375,418.24032156)
\lineto(730.35703125,418.24032156)
\lineto(730.35703125,428.54344656)
\lineto(732.7359375,428.54344656)
\lineto(732.7359375,427.44422781)
\curveto(733.0640625,427.87899291)(733.4291013,428.20301661)(733.83515625,428.41219656)
\curveto(734.23710885,428.61727521)(734.69648385,428.72391531)(735.21328125,428.72391531)
\curveto(736.13203125,428.72391531)(736.88671875,428.35477521)(737.47734375,427.62469656)
\curveto(738.06796875,426.89051661)(738.36328125,425.95125906)(738.36328125,424.80282156)
\curveto(738.36328125,423.65438406)(738.06796875,422.71102521)(737.47734375,421.98094656)
\curveto(736.88671875,421.24676661)(736.13203125,420.88172781)(735.21328125,420.88172781)
\curveto(734.69648385,420.88172781)(734.23710885,420.98836844)(733.83515625,421.19344656)
\curveto(733.4291013,421.39852521)(733.0640625,421.72254786)(732.7359375,422.16141531)
\closepath
\moveto(734.32734375,426.98485281)
\curveto(733.81054635,426.98485281)(733.41679635,426.79618041)(733.14609375,426.42704031)
\curveto(732.8712888,426.05379786)(732.7359375,425.51239161)(732.7359375,424.80282156)
\curveto(732.7359375,424.10145411)(732.8712888,423.56004786)(733.14609375,423.17860281)
\curveto(733.41679635,422.80536036)(733.81054635,422.62079031)(734.32734375,422.62079031)
\curveto(734.8400388,422.62079031)(735.2296875,422.80536036)(735.4921875,423.17860281)
\curveto(735.76289115,423.54774291)(735.90234375,424.08915021)(735.90234375,424.80282156)
\curveto(735.90234375,425.51239161)(735.76289115,426.05379786)(735.4921875,426.42704031)
\curveto(735.2296875,426.79618041)(734.8400388,426.98485281)(734.32734375,426.98485281)
\closepath
\moveto(734.32734375,426.98485281)
}
}
{
\newrgbcolor{curcolor}{0 0 0}
\pscustom[linestyle=none,fillstyle=solid,fillcolor=curcolor]
{
\newpath
\moveto(743.6625,427.01766531)
\curveto(743.1375,427.01766531)(742.73144505,426.82489161)(742.4484375,426.44344656)
\curveto(742.17363255,426.07020411)(742.03828125,425.52469656)(742.03828125,424.80282156)
\curveto(742.03828125,424.08915021)(742.17363255,423.54364161)(742.4484375,423.16219656)
\curveto(742.73144505,422.77665021)(743.1375,422.58797781)(743.6625,422.58797781)
\curveto(744.1875,422.58797781)(744.5853513,422.77665021)(744.86015625,423.16219656)
\curveto(745.1308599,423.54364161)(745.2703125,424.08915021)(745.2703125,424.80282156)
\curveto(745.2703125,425.52469656)(745.1308599,426.07020411)(744.86015625,426.44344656)
\curveto(744.5853513,426.82489161)(744.1875,427.01766531)(743.6625,427.01766531)
\closepath
\moveto(743.6625,428.72391531)
\curveto(744.9503901,428.72391531)(745.959375,428.37118041)(746.68125,427.67391531)
\curveto(747.403125,426.98485281)(747.7640625,426.02508771)(747.7640625,424.80282156)
\curveto(747.7640625,423.57645411)(747.403125,422.61258771)(746.68125,421.91532156)
\curveto(745.959375,421.22625906)(744.9503901,420.88172781)(743.6625,420.88172781)
\curveto(742.3828125,420.88172781)(741.3738276,421.22625906)(740.64375,421.91532156)
\curveto(739.921875,422.61258771)(739.5609375,423.57645411)(739.5609375,424.80282156)
\curveto(739.5609375,426.02508771)(739.921875,426.98485281)(740.64375,427.67391531)
\curveto(741.3738276,428.37118041)(742.3828125,428.72391531)(743.6625,428.72391531)
\closepath
\moveto(743.6625,428.72391531)
}
}
{
\newrgbcolor{curcolor}{0 0 0}
\pscustom[linestyle=none,fillstyle=solid,fillcolor=curcolor]
{
\newpath
\moveto(756.98847656,425.62313406)
\lineto(756.98847656,421.07860281)
\lineto(754.59316406,421.07860281)
\lineto(754.59316406,424.55672781)
\curveto(754.59316406,425.20067286)(754.57675781,425.64364161)(754.54394531,425.88563406)
\curveto(754.51933541,426.13582911)(754.47011771,426.31629786)(754.39628906,426.42704031)
\curveto(754.30605521,426.57879786)(754.18300781,426.69774291)(754.01894531,426.78797781)
\curveto(753.86308541,426.87411036)(753.68261771,426.91922781)(753.47753906,426.91922781)
\curveto(752.97304661,426.91922781)(752.57929661,426.72645411)(752.29628906,426.34500906)
\curveto(752.00918021,425.95946271)(751.86972656,425.43036036)(751.86972656,424.75360281)
\lineto(751.86972656,421.07860281)
\lineto(749.49082031,421.07860281)
\lineto(749.49082031,428.54344656)
\lineto(751.86972656,428.54344656)
\lineto(751.86972656,427.44422781)
\curveto(752.23066406,427.87899291)(752.61210911,428.20301661)(753.01816406,428.41219656)
\curveto(753.42011771,428.61727521)(753.87128906,428.72391531)(754.36347656,428.72391531)
\curveto(755.22480521,428.72391531)(755.87695286,428.45321271)(756.31582031,427.92000906)
\curveto(756.76289036,427.39500906)(756.98847656,426.62801661)(756.98847656,425.62313406)
\closepath
\moveto(756.98847656,425.62313406)
}
}
{
\newrgbcolor{curcolor}{0 0 0}
\pscustom[linestyle=none,fillstyle=solid,fillcolor=curcolor]
{
\newpath
\moveto(764.28105469,427.44422781)
\lineto(764.28105469,431.44735281)
\lineto(766.69277344,431.44735281)
\lineto(766.69277344,421.07860281)
\lineto(764.28105469,421.07860281)
\lineto(764.28105469,422.16141531)
\curveto(763.95292969,421.72254786)(763.59199219,421.39852521)(763.19824219,421.19344656)
\curveto(762.80449219,420.98836844)(762.34921849,420.88172781)(761.83652344,420.88172781)
\curveto(760.90546849,420.88172781)(760.14667969,421.24676661)(759.55605469,421.98094656)
\curveto(758.96542969,422.71102521)(758.67011719,423.65438406)(758.67011719,424.80282156)
\curveto(758.67011719,425.95125906)(758.96542969,426.89051661)(759.55605469,427.62469656)
\curveto(760.14667969,428.35477521)(760.90546849,428.72391531)(761.83652344,428.72391531)
\curveto(762.34921849,428.72391531)(762.80449219,428.61727521)(763.19824219,428.41219656)
\curveto(763.59199219,428.20301661)(763.95292969,427.87899291)(764.28105469,427.44422781)
\closepath
\moveto(762.72246094,422.62079031)
\curveto(763.22285104,422.62079031)(763.60839844,422.80536036)(763.87089844,423.17860281)
\curveto(764.14160104,423.54774291)(764.28105469,424.08915021)(764.28105469,424.80282156)
\curveto(764.28105469,425.51239161)(764.14160104,426.05379786)(763.87089844,426.42704031)
\curveto(763.60839844,426.79618041)(763.22285104,426.98485281)(762.72246094,426.98485281)
\curveto(762.20566354,426.98485281)(761.81191354,426.79618041)(761.54121094,426.42704031)
\curveto(761.27871094,426.05379786)(761.14746094,425.51239161)(761.14746094,424.80282156)
\curveto(761.14746094,424.08915021)(761.27871094,423.54774291)(761.54121094,423.17860281)
\curveto(761.81191354,422.80536036)(762.20566354,422.62079031)(762.72246094,422.62079031)
\closepath
\moveto(762.72246094,422.62079031)
}
}
{
\newrgbcolor{curcolor}{0 0 0}
\pscustom[linestyle=none,fillstyle=solid,fillcolor=curcolor]
{
\newpath
\moveto(776.42167969,424.83563406)
\lineto(776.42167969,424.14657156)
\lineto(770.84355469,424.14657156)
\curveto(770.89687474,423.58875906)(771.09785209,423.16629786)(771.45058594,422.88329031)
\curveto(771.79921849,422.60848536)(772.28730469,422.47313406)(772.91074219,422.47313406)
\curveto(773.40292969,422.47313406)(773.91152344,422.54286036)(774.43652344,422.68641531)
\curveto(774.96972709,422.83817286)(775.51933594,423.06786036)(776.07714844,423.37547781)
\lineto(776.07714844,421.53797781)
\curveto(775.50703099,421.31649291)(774.93691459,421.15243094)(774.37089844,421.04579031)
\curveto(773.81308594,420.93915021)(773.24706979,420.88172781)(772.68105469,420.88172781)
\curveto(771.33574219,420.88172781)(770.28574219,421.22625906)(769.53105469,421.91532156)
\curveto(768.78456979,422.60438406)(768.41542969,423.56415021)(768.41542969,424.80282156)
\curveto(768.41542969,426.01688406)(768.78046849,426.97254786)(769.51464844,427.67391531)
\curveto(770.24472709,428.37118041)(771.25781224,428.72391531)(772.54980469,428.72391531)
\curveto(773.71874974,428.72391531)(774.65390599,428.36707911)(775.35527344,427.65750906)
\curveto(776.06484349,426.95614161)(776.42167969,426.01688406)(776.42167969,424.83563406)
\closepath
\moveto(773.97714844,425.62313406)
\curveto(773.97714844,426.08250906)(773.83769479,426.44754786)(773.56699219,426.72235281)
\curveto(773.30449219,426.99305541)(772.95996094,427.13250906)(772.53339844,427.13250906)
\curveto(772.07402344,427.13250906)(771.69667969,427.00125906)(771.40136719,426.73875906)
\curveto(771.11425729,426.48446271)(770.94199219,426.11532156)(770.87636719,425.62313406)
\closepath
\moveto(773.97714844,425.62313406)
}
}
{
\newrgbcolor{curcolor}{0 0 0}
\pscustom[linestyle=none,fillstyle=solid,fillcolor=curcolor]
{
\newpath
\moveto(783.76757812,426.50907156)
\curveto(783.55839817,426.60750906)(783.34921822,426.67723536)(783.14414062,426.72235281)
\curveto(782.93496067,426.77567286)(782.72578072,426.80438406)(782.52070312,426.80438406)
\curveto(781.90546822,426.80438406)(781.42968803,426.60750906)(781.09335937,426.21375906)
\curveto(780.76523437,425.82000906)(780.60117187,425.25399291)(780.60117187,424.52391531)
\lineto(780.60117187,421.07860281)
\lineto(778.22226562,421.07860281)
\lineto(778.22226562,428.54344656)
\lineto(780.60117187,428.54344656)
\lineto(780.60117187,427.31297781)
\curveto(780.90468802,427.80516531)(781.25742187,428.15790021)(781.65117187,428.37938406)
\curveto(782.05312552,428.60907156)(782.53710938,428.72391531)(783.09492187,428.72391531)
\curveto(783.16875052,428.72391531)(783.25078073,428.71571271)(783.34101562,428.70750906)
\curveto(783.43945312,428.70750906)(783.57480443,428.69520411)(783.75117187,428.67469656)
\closepath
\moveto(783.76757812,426.50907156)
}
}
{
\newrgbcolor{curcolor}{0 0 0}
\pscustom[linewidth=2.10000003,linecolor=curcolor]
{
\newpath
\moveto(143.91972,436.25438275)
\curveto(143.91972,417.92451175)(114.199806,403.06453375)(77.540043,403.06453375)
\curveto(40.88028,403.06453375)(11.160345,417.92451175)(11.160345,436.25438275)
\curveto(11.160345,454.58427475)(40.88028,469.44423175)(77.540043,469.44423175)
\curveto(114.199806,469.44423175)(143.91972,454.58427475)(143.91972,436.25438275)
}
}
{
\newrgbcolor{curcolor}{0 0 0}
\pscustom[linestyle=none,fillstyle=solid,fillcolor=curcolor]
{
\newpath
\moveto(64.30429688,442.11961844)
\lineto(67.10976562,442.11961844)
\lineto(69.37382813,438.55946219)
\lineto(71.63789062,442.11961844)
\lineto(74.45976563,442.11961844)
\lineto(70.65351562,436.36102469)
\lineto(70.65351562,432.16102469)
\lineto(68.09414062,432.16102469)
\lineto(68.09414062,436.36102469)
\closepath
\moveto(64.30429688,442.11961844)
}
}
{
\newrgbcolor{curcolor}{0 0 0}
\pscustom[linestyle=none,fillstyle=solid,fillcolor=curcolor]
{
\newpath
\moveto(81.87128906,435.91805594)
\lineto(81.87128906,435.22899344)
\lineto(76.29316406,435.22899344)
\curveto(76.34648411,434.67118094)(76.54746146,434.24871974)(76.90019531,433.96571219)
\curveto(77.24882786,433.69090724)(77.73691406,433.55555594)(78.36035156,433.55555594)
\curveto(78.85253906,433.55555594)(79.36113281,433.62528224)(79.88613281,433.76883719)
\curveto(80.41933646,433.92059474)(80.96894531,434.15028224)(81.52675781,434.45789969)
\lineto(81.52675781,432.62039969)
\curveto(80.95664036,432.39891479)(80.38652396,432.23485281)(79.82050781,432.12821219)
\curveto(79.26269531,432.02157209)(78.69667916,431.96414969)(78.13066406,431.96414969)
\curveto(76.78535156,431.96414969)(75.73535156,432.30868094)(74.98066406,432.99774344)
\curveto(74.23417916,433.68680594)(73.86503906,434.64657209)(73.86503906,435.88524344)
\curveto(73.86503906,437.09930594)(74.23007786,438.05496974)(74.96425781,438.75633719)
\curveto(75.69433646,439.45360229)(76.70742161,439.80633719)(77.99941406,439.80633719)
\curveto(79.16835911,439.80633719)(80.10351536,439.44950099)(80.80488281,438.73993094)
\curveto(81.51445286,438.03856349)(81.87128906,437.09930594)(81.87128906,435.91805594)
\closepath
\moveto(79.42675781,436.70555594)
\curveto(79.42675781,437.16493094)(79.28730416,437.52996974)(79.01660156,437.80477469)
\curveto(78.75410156,438.07547729)(78.40957031,438.21493094)(77.98300781,438.21493094)
\curveto(77.52363281,438.21493094)(77.14628906,438.08368094)(76.85097656,437.82118094)
\curveto(76.56386666,437.56688459)(76.39160156,437.19774344)(76.32597656,436.70555594)
\closepath
\moveto(79.42675781,436.70555594)
}
}
{
\newrgbcolor{curcolor}{0 0 0}
\pscustom[linestyle=none,fillstyle=solid,fillcolor=curcolor]
{
\newpath
\moveto(89.49609375,439.39618094)
\lineto(89.49609375,437.57508719)
\curveto(88.9916013,437.79246974)(88.4994138,437.95243094)(88.01953125,438.05086844)
\curveto(87.5478513,438.15750959)(87.10078125,438.21493094)(86.67421875,438.21493094)
\curveto(86.2230474,438.21493094)(85.88671875,438.15750959)(85.65703125,438.05086844)
\curveto(85.4355474,437.94012599)(85.32890625,437.76375959)(85.32890625,437.52586844)
\curveto(85.32890625,437.33719709)(85.4109375,437.18953979)(85.575,437.08289969)
\curveto(85.7472651,436.98446219)(86.05078125,436.90653224)(86.47734375,436.85321219)
\lineto(86.90390625,436.80399344)
\curveto(88.1261724,436.64813459)(88.9464849,436.39383719)(89.36484375,436.03289969)
\curveto(89.79140625,435.68016479)(90.0046875,435.12235229)(90.0046875,434.35946219)
\curveto(90.0046875,433.55965724)(89.709375,432.95672729)(89.11875,432.55477469)
\curveto(88.5363276,432.16102469)(87.66269505,431.96414969)(86.49375,431.96414969)
\curveto(86.0015625,431.96414969)(85.49296875,432.00516479)(84.96796875,432.07899344)
\curveto(84.44296875,432.15692312)(83.9015625,432.27586844)(83.34375,432.43993094)
\lineto(83.34375,434.24461844)
\curveto(83.8113276,434.01493094)(84.2994138,433.83856349)(84.80390625,433.71961844)
\curveto(85.3042974,433.60887599)(85.8128901,433.55555594)(86.3296875,433.55555594)
\curveto(86.7972651,433.55555594)(87.15,433.61297729)(87.3796875,433.73602469)
\curveto(87.6175776,433.86727469)(87.740625,434.06414969)(87.740625,434.32664969)
\curveto(87.740625,434.53172729)(87.65859375,434.68758719)(87.49453125,434.78602469)
\curveto(87.33046875,434.89266479)(87.00644505,434.98289969)(86.5265625,435.04852469)
\lineto(86.1,435.09774344)
\curveto(85.03769505,435.22899344)(84.2953125,435.47508719)(83.86875,435.83602469)
\curveto(83.4421875,436.19696219)(83.22890625,436.74246974)(83.22890625,437.47664969)
\curveto(83.22890625,438.26414969)(83.49550755,438.84657209)(84.0328125,439.23211844)
\curveto(84.57832005,439.61356349)(85.4150388,439.80633719)(86.54296875,439.80633719)
\curveto(86.9777349,439.80633719)(87.4371099,439.77352469)(87.92109375,439.70789969)
\curveto(88.4009763,439.64227469)(88.9259763,439.53563459)(89.49609375,439.39618094)
\closepath
\moveto(89.49609375,439.39618094)
}
}
{
\newrgbcolor{curcolor}{0 0 0}
\pscustom[linewidth=2.10000003,linecolor=curcolor]
{
\newpath
\moveto(364.514052,434.73270175)
\curveto(364.514052,416.40283075)(334.794138,401.54285275)(298.134375,401.54285275)
\curveto(261.474612,401.54285275)(231.754677,416.40283075)(231.754677,434.73270175)
\curveto(231.754677,453.06259375)(261.474612,467.92255075)(298.134375,467.92255075)
\curveto(334.794138,467.92255075)(364.514052,453.06259375)(364.514052,434.73270175)
}
}
{
\newrgbcolor{curcolor}{0 0 0}
\pscustom[linestyle=none,fillstyle=solid,fillcolor=curcolor]
{
\newpath
\moveto(288.98378906,440.59383719)
\lineto(291.85488281,440.59383719)
\lineto(295.48066406,433.76883719)
\lineto(295.48066406,440.59383719)
\lineto(297.90878906,440.59383719)
\lineto(297.90878906,430.63524344)
\lineto(295.03769531,430.63524344)
\lineto(291.42832031,437.46024344)
\lineto(291.42832031,430.63524344)
\lineto(288.98378906,430.63524344)
\closepath
\moveto(288.98378906,440.59383719)
}
}
{
\newrgbcolor{curcolor}{0 0 0}
\pscustom[linestyle=none,fillstyle=solid,fillcolor=curcolor]
{
\newpath
\moveto(303.85195313,436.57430594)
\curveto(303.32695313,436.57430594)(302.92089818,436.38153224)(302.63789063,436.00008719)
\curveto(302.36308568,435.62684474)(302.22773438,435.08133719)(302.22773438,434.35946219)
\curveto(302.22773438,433.64579084)(302.36308568,433.10028224)(302.63789063,432.71883719)
\curveto(302.92089818,432.33329084)(303.32695313,432.14461844)(303.85195313,432.14461844)
\curveto(304.37695313,432.14461844)(304.77480443,432.33329084)(305.04960938,432.71883719)
\curveto(305.32031303,433.10028224)(305.45976563,433.64579084)(305.45976563,434.35946219)
\curveto(305.45976563,435.08133719)(305.32031303,435.62684474)(305.04960938,436.00008719)
\curveto(304.77480443,436.38153224)(304.37695313,436.57430594)(303.85195313,436.57430594)
\closepath
\moveto(303.85195313,438.28055594)
\curveto(305.13984323,438.28055594)(306.14882813,437.92782104)(306.87070313,437.23055594)
\curveto(307.59257813,436.54149344)(307.95351563,435.58172834)(307.95351563,434.35946219)
\curveto(307.95351563,433.13309474)(307.59257813,432.16922834)(306.87070313,431.47196219)
\curveto(306.14882813,430.78289969)(305.13984323,430.43836844)(303.85195313,430.43836844)
\curveto(302.57226563,430.43836844)(301.56328073,430.78289969)(300.83320313,431.47196219)
\curveto(300.11132813,432.16922834)(299.75039063,433.13309474)(299.75039063,434.35946219)
\curveto(299.75039063,435.58172834)(300.11132813,436.54149344)(300.83320313,437.23055594)
\curveto(301.56328073,437.92782104)(302.57226563,438.28055594)(303.85195313,438.28055594)
\closepath
\moveto(303.85195313,438.28055594)
}
}
{
\newrgbcolor{curcolor}{0 0 0}
\pscustom[linewidth=2.10000003,linecolor=curcolor]
{
\newpath
\moveto(585.104289,432.16102075)
\curveto(585.104289,413.83114975)(555.384375,398.97117175)(518.724612,398.97117175)
\curveto(482.064849,398.97117175)(452.344914,413.83114975)(452.344914,432.16102075)
\curveto(452.344914,450.49091275)(482.064849,465.35086975)(518.724612,465.35086975)
\curveto(555.384375,465.35086975)(585.104289,450.49091275)(585.104289,432.16102075)
}
}
{
\newrgbcolor{curcolor}{0 0 0}
\pscustom[linestyle=none,fillstyle=solid,fillcolor=curcolor]
{
\newpath
\moveto(494.62382812,438.0221575)
\lineto(497.88867188,438.0221575)
\lineto(500.16914063,432.69012625)
\lineto(502.44960938,438.0221575)
\lineto(505.69804687,438.0221575)
\lineto(505.69804687,428.06356375)
\lineto(503.26992188,428.06356375)
\lineto(503.26992188,435.34793875)
\lineto(500.98945313,429.983095)
\lineto(499.34882813,429.983095)
\lineto(497.06835938,435.34793875)
\lineto(497.06835938,428.06356375)
\lineto(494.62382812,428.06356375)
\closepath
\moveto(494.62382812,438.0221575)
}
}
{
\newrgbcolor{curcolor}{0 0 0}
\pscustom[linestyle=none,fillstyle=solid,fillcolor=curcolor]
{
\newpath
\moveto(511.4484375,431.426845)
\curveto(510.94394505,431.426845)(510.5666013,431.33661115)(510.31640625,431.164345)
\curveto(510.0621099,431.0002825)(509.9390625,430.75418875)(509.9390625,430.42606375)
\curveto(509.9390625,430.1184463)(510.0375,429.87645385)(510.234375,429.70418875)
\curveto(510.4394526,429.54012625)(510.7265625,429.458095)(511.0875,429.458095)
\curveto(511.53457005,429.458095)(511.9119138,429.61395385)(512.21953125,429.93387625)
\curveto(512.5230474,430.26200125)(512.67890625,430.66395385)(512.67890625,431.14793875)
\lineto(512.67890625,431.426845)
\closepath
\moveto(515.090625,432.32918875)
\lineto(515.090625,428.06356375)
\lineto(512.67890625,428.06356375)
\lineto(512.67890625,429.1627825)
\curveto(512.3589849,428.71161115)(511.9980474,428.38348615)(511.59609375,428.1784075)
\curveto(511.20234375,427.97332937)(510.7183599,427.86668875)(510.15234375,427.86668875)
\curveto(509.3730474,427.86668875)(508.74550755,428.08817312)(508.265625,428.539345)
\curveto(507.7816401,428.99872)(507.54375,429.589345)(507.54375,430.31122)
\curveto(507.54375,431.18485255)(507.8431638,431.8246963)(508.44609375,432.23075125)
\curveto(509.0449224,432.63270385)(509.99238255,432.8377825)(511.284375,432.8377825)
\lineto(512.67890625,432.8377825)
\lineto(512.67890625,433.0346575)
\curveto(512.67890625,433.40379865)(512.53125,433.67860255)(512.2359375,433.85497)
\curveto(511.940625,434.02723615)(511.4730474,434.11747)(510.84140625,434.11747)
\curveto(510.3369138,434.11747)(509.86113255,434.06825125)(509.4140625,433.96981375)
\curveto(508.97519505,433.87137625)(508.5691401,433.71551635)(508.2,433.51043875)
\lineto(508.2,435.3315325)
\curveto(508.7003901,435.45047755)(509.2089849,435.54481375)(509.72578125,435.61043875)
\curveto(510.25078125,435.67606375)(510.7675776,435.70887625)(511.284375,435.70887625)
\curveto(512.6296875,435.70887625)(513.59765625,435.43817365)(514.18828125,434.90497)
\curveto(514.7871099,434.37997)(515.090625,433.51864135)(515.090625,432.32918875)
\closepath
\moveto(515.090625,432.32918875)
}
}
{
\newrgbcolor{curcolor}{0 0 0}
\pscustom[linestyle=none,fillstyle=solid,fillcolor=curcolor]
{
\newpath
\moveto(516.32519531,435.5284075)
\lineto(518.72050781,435.5284075)
\lineto(520.72207031,430.45887625)
\lineto(522.42832031,435.5284075)
\lineto(524.80722656,435.5284075)
\lineto(521.67363281,427.358095)
\curveto(521.35371041,426.52957885)(520.98457031,425.9471575)(520.55800781,425.6190325)
\curveto(520.13964791,425.28270385)(519.59003906,425.11043875)(518.90097656,425.11043875)
\lineto(517.52285156,425.11043875)
\lineto(517.52285156,426.68543875)
\lineto(518.27753906,426.68543875)
\curveto(518.67949271,426.68543875)(518.97070286,426.75106375)(519.14707031,426.88231375)
\curveto(519.33164036,427.00536115)(519.47929661,427.23504865)(519.59003906,427.57137625)
\lineto(519.65566406,427.7846575)
\closepath
\moveto(516.32519531,435.5284075)
}
}
{
\newrgbcolor{curcolor}{0 0 0}
\pscustom[linestyle=none,fillstyle=solid,fillcolor=curcolor]
{
\newpath
\moveto(530.18027344,429.60575125)
\curveto(530.69296849,429.60575125)(531.08261719,429.7903213)(531.34511719,430.16356375)
\curveto(531.61582084,430.53270385)(531.75527344,431.07411115)(531.75527344,431.7877825)
\curveto(531.75527344,432.49735255)(531.61582084,433.0387588)(531.34511719,433.41200125)
\curveto(531.08261719,433.78114135)(530.69296849,433.96981375)(530.18027344,433.96981375)
\curveto(529.66347604,433.96981375)(529.26972604,433.78114135)(528.99902344,433.41200125)
\curveto(528.72421849,433.0387588)(528.58886719,432.49735255)(528.58886719,431.7877825)
\curveto(528.58886719,431.08641505)(528.72421849,430.5450088)(528.99902344,430.16356375)
\curveto(529.26972604,429.7903213)(529.66347604,429.60575125)(530.18027344,429.60575125)
\closepath
\moveto(528.58886719,434.42918875)
\curveto(528.91699219,434.86395385)(529.28203099,435.18797755)(529.68808594,435.3971575)
\curveto(530.09003854,435.60223615)(530.54941354,435.70887625)(531.06621094,435.70887625)
\curveto(531.98496094,435.70887625)(532.73964844,435.33973615)(533.33027344,434.6096575)
\curveto(533.92089844,433.87547755)(534.21621094,432.93622)(534.21621094,431.7877825)
\curveto(534.21621094,430.639345)(533.92089844,429.69598615)(533.33027344,428.9659075)
\curveto(532.73964844,428.23172755)(531.98496094,427.86668875)(531.06621094,427.86668875)
\curveto(530.54941354,427.86668875)(530.09003854,427.97332937)(529.68808594,428.1784075)
\curveto(529.28203099,428.38348615)(528.91699219,428.7075088)(528.58886719,429.14637625)
\lineto(528.58886719,428.06356375)
\lineto(526.20996094,428.06356375)
\lineto(526.20996094,438.43231375)
\lineto(528.58886719,438.43231375)
\closepath
\moveto(528.58886719,434.42918875)
}
}
{
\newrgbcolor{curcolor}{0 0 0}
\pscustom[linestyle=none,fillstyle=solid,fillcolor=curcolor]
{
\newpath
\moveto(543.42011719,431.820595)
\lineto(543.42011719,431.1315325)
\lineto(537.84199219,431.1315325)
\curveto(537.89531224,430.57372)(538.09628959,430.1512588)(538.44902344,429.86825125)
\curveto(538.79765599,429.5934463)(539.28574219,429.458095)(539.90917969,429.458095)
\curveto(540.40136719,429.458095)(540.90996094,429.5278213)(541.43496094,429.67137625)
\curveto(541.96816459,429.8231338)(542.51777344,430.0528213)(543.07558594,430.36043875)
\lineto(543.07558594,428.52293875)
\curveto(542.50546849,428.30145385)(541.93535209,428.13739187)(541.36933594,428.03075125)
\curveto(540.81152344,427.92411115)(540.24550729,427.86668875)(539.67949219,427.86668875)
\curveto(538.33417969,427.86668875)(537.28417969,428.21122)(536.52949219,428.9002825)
\curveto(535.78300729,429.589345)(535.41386719,430.54911115)(535.41386719,431.7877825)
\curveto(535.41386719,433.001845)(535.77890599,433.9575088)(536.51308594,434.65887625)
\curveto(537.24316459,435.35614135)(538.25624974,435.70887625)(539.54824219,435.70887625)
\curveto(540.71718724,435.70887625)(541.65234349,435.35204005)(542.35371094,434.64247)
\curveto(543.06328099,433.94110255)(543.42011719,433.001845)(543.42011719,431.820595)
\closepath
\moveto(540.97558594,432.608095)
\curveto(540.97558594,433.06747)(540.83613229,433.4325088)(540.56542969,433.70731375)
\curveto(540.30292969,433.97801635)(539.95839844,434.11747)(539.53183594,434.11747)
\curveto(539.07246094,434.11747)(538.69511719,433.98622)(538.39980469,433.72372)
\curveto(538.11269479,433.46942365)(537.94042969,433.1002825)(537.87480469,432.608095)
\closepath
\moveto(540.97558594,432.608095)
}
}
{
\newrgbcolor{curcolor}{0 0 0}
\pscustom[linewidth=2.10000003,linecolor=curcolor]
{
\newpath
\moveto(331.586724,548.02606375)
\curveto(329.933814,477.67606375)(77.909181,524.92606375)(77.609763,479.66533075)
}
}
{
\newrgbcolor{curcolor}{0 0 0}
\pscustom[linestyle=none,fillstyle=solid,fillcolor=curcolor]
{
\newpath
\moveto(77.556444,471.79033075)
\lineto(72.376164,482.32722775)
\lineto(77.609763,479.66533075)
\lineto(82.876164,482.25750775)
\closepath
\moveto(77.556444,471.79033075)
}
}
{
\newrgbcolor{curcolor}{0 0 0}
\pscustom[linewidth=2.10000003,linecolor=curcolor]
{
\newpath
\moveto(77.556444,471.79033075)
\lineto(72.376164,482.32722775)
\lineto(77.609763,479.66533075)
\lineto(82.876164,482.25750775)
\closepath
\moveto(77.556444,471.79033075)
}
}
{
\newrgbcolor{curcolor}{0 0 0}
\pscustom[linewidth=2.10000003,linecolor=curcolor]
{
\newpath
\moveto(367.221099,548.02606375)
\curveto(367.721529,510.22606375)(300.509181,497.62606375)(298.946487,478.11082675)
}
}
{
\newrgbcolor{curcolor}{0 0 0}
\pscustom[linestyle=none,fillstyle=solid,fillcolor=curcolor]
{
\newpath
\moveto(298.318944,470.26043875)
\lineto(293.922069,481.14597775)
\lineto(298.946487,478.11082675)
\lineto(304.389267,480.30927475)
\closepath
\moveto(298.318944,470.26043875)
}
}
{
\newrgbcolor{curcolor}{0 0 0}
\pscustom[linewidth=2.10000003,linecolor=curcolor]
{
\newpath
\moveto(298.318944,470.26043875)
\lineto(293.922069,481.14597775)
\lineto(298.946487,478.11082675)
\lineto(304.389267,480.30927475)
\closepath
\moveto(298.318944,470.26043875)
}
}
{
\newrgbcolor{curcolor}{0 0 0}
\pscustom[linewidth=2.10000003,linecolor=curcolor]
{
\newpath
\moveto(438.489849,548.02606375)
\curveto(440.146959,495.52606375)(517.859181,502.87606375)(518.490819,475.56785275)
}
}
{
\newrgbcolor{curcolor}{0 0 0}
\pscustom[linestyle=none,fillstyle=solid,fillcolor=curcolor]
{
\newpath
\moveto(518.671293,467.69696875)
\lineto(513.179289,478.07390875)
\lineto(518.490819,475.56785275)
\lineto(523.675194,478.31591275)
\closepath
\moveto(518.671293,467.69696875)
}
}
{
\newrgbcolor{curcolor}{0 0 0}
\pscustom[linewidth=2.10000003,linecolor=curcolor]
{
\newpath
\moveto(518.671293,467.69696875)
\lineto(513.179289,478.07390875)
\lineto(518.490819,475.56785275)
\lineto(523.675194,478.31591275)
\closepath
\moveto(518.671293,467.69696875)
}
}
{
\newrgbcolor{curcolor}{0 0 0}
\pscustom[linewidth=2.10000003,linecolor=curcolor]
{
\newpath
\moveto(474.124224,548.02606375)
\curveto(476.884569,478.72606375)(740.459181,521.77606375)(742.15722,478.59890875)
}
}
{
\newrgbcolor{curcolor}{0 0 0}
\pscustom[linestyle=none,fillstyle=solid,fillcolor=curcolor]
{
\newpath
\moveto(742.464849,470.73211975)
\lineto(736.808793,481.01884375)
\lineto(742.15722,478.59890875)
\lineto(747.300582,481.42899475)
\closepath
\moveto(742.464849,470.73211975)
}
}
{
\newrgbcolor{curcolor}{0 0 0}
\pscustom[linewidth=2.10000003,linecolor=curcolor]
{
\newpath
\moveto(742.464849,470.73211975)
\lineto(736.808793,481.01884375)
\lineto(742.15722,478.59890875)
\lineto(747.300582,481.42899475)
\closepath
\moveto(742.464849,470.73211975)
}
}
{
\newrgbcolor{curcolor}{0 0 0}
\pscustom[linewidth=2.10000003,linecolor=curcolor]
{
\newpath
\moveto(1.05,305.65653775)
\lineto(151.569138,305.65653775)
\lineto(151.569138,265.75653775)
\lineto(1.05,265.75653775)
\closepath
\moveto(1.05,305.65653775)
}
}
{
\newrgbcolor{curcolor}{0 0 0}
\pscustom[linestyle=none,fillstyle=solid,fillcolor=curcolor]
{
\newpath
\moveto(28.09570312,281.43270437)
\lineto(27.89882813,281.43270437)
\curveto(26.25820313,281.43270437)(24.97851562,281.88387572)(24.05976563,282.79442312)
\curveto(23.14921822,283.70086817)(22.69804687,284.96414942)(22.69804687,286.58426687)
\curveto(22.69804687,288.19207937)(23.14921822,289.45536062)(24.05976563,290.37411062)
\curveto(24.97851562,291.29286062)(26.22949192,291.75223562)(27.81679688,291.75223562)
\curveto(29.43281303,291.75223562)(30.69199193,291.29696192)(31.59023438,290.39051687)
\curveto(32.48437553,289.47997052)(32.93554688,288.21258692)(32.93554688,286.58426687)
\curveto(32.93554688,285.45633692)(32.69355443,284.48836817)(32.21367188,283.68036062)
\curveto(31.74199193,282.86825072)(31.05703073,282.26122052)(30.16289063,281.85926687)
\lineto(32.16445313,279.62801687)
\lineto(29.71992188,279.62801687)
\closepath
\moveto(27.81679688,289.88192312)
\curveto(27.03750052,289.88192312)(26.43046822,289.59071192)(25.99570312,289.01239187)
\curveto(25.56914062,288.44227442)(25.35585938,287.63426687)(25.35585938,286.58426687)
\curveto(25.35585938,285.50965802)(25.56093802,284.68934552)(25.97929687,284.12332937)
\curveto(26.40585937,283.56551687)(27.01699192,283.28661062)(27.81679688,283.28661062)
\curveto(28.60429687,283.28661062)(29.21132813,283.56961817)(29.63789063,284.13973562)
\curveto(30.07265573,284.71805567)(30.29414063,285.53426687)(30.29414063,286.58426687)
\curveto(30.29414063,287.63426687)(30.07265573,288.44227442)(29.63789063,289.01239187)
\curveto(29.21132813,289.59071192)(28.60429687,289.88192312)(27.81679688,289.88192312)
\closepath
\moveto(27.81679688,289.88192312)
}
}
{
\newrgbcolor{curcolor}{0 0 0}
\pscustom[linestyle=none,fillstyle=solid,fillcolor=curcolor]
{
\newpath
\moveto(34.69921875,284.51707937)
\lineto(34.69921875,289.07801687)
\lineto(37.09453125,289.07801687)
\lineto(37.09453125,288.33973562)
\curveto(37.09453125,287.93368067)(37.0863276,287.42508692)(37.078125,286.81395437)
\lineto(37.078125,285.58348562)
\curveto(37.078125,284.98055567)(37.09453125,284.54989187)(37.12734375,284.28739187)
\curveto(37.16015625,284.02489187)(37.2134763,283.83211817)(37.29140625,283.71317312)
\curveto(37.38984375,283.55731322)(37.5128901,283.43836817)(37.66875,283.35223562)
\curveto(37.8328125,283.26200072)(38.01738255,283.22098562)(38.2265625,283.22098562)
\curveto(38.71875,283.22098562)(39.1042974,283.40965802)(39.39140625,283.79520437)
\curveto(39.6744138,284.17664942)(39.81796875,284.70575072)(39.81796875,285.38661062)
\lineto(39.81796875,289.07801687)
\lineto(42.196875,289.07801687)
\lineto(42.196875,281.61317312)
\lineto(39.81796875,281.61317312)
\lineto(39.81796875,282.69598562)
\curveto(39.45703125,282.25711817)(39.0714849,281.93309552)(38.66953125,281.72801687)
\curveto(38.27578125,281.52293875)(37.8328125,281.41629812)(37.340625,281.41629812)
\curveto(36.47519505,281.41629812)(35.81894505,281.67879812)(35.371875,282.20379812)
\curveto(34.9207026,282.73700072)(34.69921875,283.50809552)(34.69921875,284.51707937)
\closepath
\moveto(34.69921875,284.51707937)
}
}
{
\newrgbcolor{curcolor}{0 0 0}
\pscustom[linestyle=none,fillstyle=solid,fillcolor=curcolor]
{
\newpath
\moveto(51.95039062,285.37020437)
\lineto(51.95039062,284.68114187)
\lineto(46.37226562,284.68114187)
\curveto(46.42558568,284.12332937)(46.62656303,283.70086817)(46.97929688,283.41786062)
\curveto(47.32792943,283.14305567)(47.81601562,283.00770437)(48.43945312,283.00770437)
\curveto(48.93164062,283.00770437)(49.44023438,283.07743067)(49.96523438,283.22098562)
\curveto(50.49843802,283.37274317)(51.04804687,283.60243067)(51.60585938,283.91004812)
\lineto(51.60585938,282.07254812)
\curveto(51.03574192,281.85106322)(50.46562553,281.68700125)(49.89960937,281.58036062)
\curveto(49.34179688,281.47372052)(48.77578073,281.41629812)(48.20976563,281.41629812)
\curveto(46.86445312,281.41629812)(45.81445312,281.76082937)(45.05976562,282.44989187)
\curveto(44.31328072,283.13895437)(43.94414063,284.09872052)(43.94414063,285.33739187)
\curveto(43.94414063,286.55145437)(44.30917942,287.50711817)(45.04335938,288.20848562)
\curveto(45.77343802,288.90575072)(46.78652317,289.25848562)(48.07851563,289.25848562)
\curveto(49.24746067,289.25848562)(50.18261692,288.90164942)(50.88398437,288.19207937)
\curveto(51.59355443,287.49071192)(51.95039062,286.55145437)(51.95039062,285.37020437)
\closepath
\moveto(49.50585938,286.15770437)
\curveto(49.50585938,286.61707937)(49.36640572,286.98211817)(49.09570312,287.25692312)
\curveto(48.83320312,287.52762572)(48.48867188,287.66707937)(48.06210937,287.66707937)
\curveto(47.60273437,287.66707937)(47.22539063,287.53582937)(46.93007813,287.27332937)
\curveto(46.64296822,287.01903302)(46.47070312,286.64989187)(46.40507813,286.15770437)
\closepath
\moveto(49.50585938,286.15770437)
}
}
{
\newrgbcolor{curcolor}{0 0 0}
\pscustom[linestyle=none,fillstyle=solid,fillcolor=curcolor]
{
\newpath
\moveto(59.57519531,288.84832937)
\lineto(59.57519531,287.02723562)
\curveto(59.07070286,287.24461817)(58.57851536,287.40457937)(58.09863281,287.50301687)
\curveto(57.62695286,287.60965802)(57.17988281,287.66707937)(56.75332031,287.66707937)
\curveto(56.30214896,287.66707937)(55.96582031,287.60965802)(55.73613281,287.50301687)
\curveto(55.51464896,287.39227442)(55.40800781,287.21590802)(55.40800781,286.97801687)
\curveto(55.40800781,286.78934552)(55.49003906,286.64168822)(55.65410156,286.53504812)
\curveto(55.82636666,286.43661062)(56.12988281,286.35868067)(56.55644531,286.30536062)
\lineto(56.98300781,286.25614187)
\curveto(58.20527396,286.10028302)(59.02558646,285.84598562)(59.44394531,285.48504812)
\curveto(59.87050781,285.13231322)(60.08378906,284.57450072)(60.08378906,283.81161062)
\curveto(60.08378906,283.01180567)(59.78847656,282.40887572)(59.19785156,282.00692312)
\curveto(58.61542916,281.61317312)(57.74179661,281.41629812)(56.57285156,281.41629812)
\curveto(56.08066406,281.41629812)(55.57207031,281.45731322)(55.04707031,281.53114187)
\curveto(54.52207031,281.60907156)(53.98066406,281.72801687)(53.42285156,281.89207937)
\lineto(53.42285156,283.69676687)
\curveto(53.89042916,283.46707937)(54.37851536,283.29071192)(54.88300781,283.17176687)
\curveto(55.38339896,283.06102442)(55.89199166,283.00770437)(56.40878906,283.00770437)
\curveto(56.87636666,283.00770437)(57.22910156,283.06512572)(57.45878906,283.18817312)
\curveto(57.69667916,283.31942312)(57.81972656,283.51629812)(57.81972656,283.77879812)
\curveto(57.81972656,283.98387572)(57.73769531,284.13973562)(57.57363281,284.23817312)
\curveto(57.40957031,284.34481322)(57.08554661,284.43504812)(56.60566406,284.50067312)
\lineto(56.17910156,284.54989187)
\curveto(55.11679661,284.68114187)(54.37441406,284.92723562)(53.94785156,285.28817312)
\curveto(53.52128906,285.64911062)(53.30800781,286.19461817)(53.30800781,286.92879812)
\curveto(53.30800781,287.71629812)(53.57460911,288.29872052)(54.11191406,288.68426687)
\curveto(54.65742161,289.06571192)(55.49414036,289.25848562)(56.62207031,289.25848562)
\curveto(57.05683646,289.25848562)(57.51621146,289.22567312)(58.00019531,289.16004812)
\curveto(58.48007786,289.09442312)(59.00507786,288.98778302)(59.57519531,288.84832937)
\closepath
\moveto(59.57519531,288.84832937)
}
}
{
\newrgbcolor{curcolor}{0 0 0}
\pscustom[linestyle=none,fillstyle=solid,fillcolor=curcolor]
{
\newpath
\moveto(64.48066406,291.19442312)
\lineto(64.48066406,289.07801687)
\lineto(66.94160156,289.07801687)
\lineto(66.94160156,287.37176687)
\lineto(64.48066406,287.37176687)
\lineto(64.48066406,284.20536062)
\curveto(64.48066406,283.85262572)(64.54628906,283.61883692)(64.67753906,283.49989187)
\curveto(64.81699271,283.37684552)(65.09179661,283.31942312)(65.49785156,283.31942312)
\lineto(66.72832031,283.31942312)
\lineto(66.72832031,281.61317312)
\lineto(64.67753906,281.61317312)
\curveto(63.73418021,281.61317312)(63.06972656,281.81004812)(62.67597656,282.20379812)
\curveto(62.28222656,282.59754812)(62.08535156,283.26200072)(62.08535156,284.20536062)
\lineto(62.08535156,287.37176687)
\lineto(60.90410156,287.37176687)
\lineto(60.90410156,289.07801687)
\lineto(62.08535156,289.07801687)
\lineto(62.08535156,291.19442312)
\closepath
\moveto(64.48066406,291.19442312)
}
}
{
\newrgbcolor{curcolor}{0 0 0}
\pscustom[linestyle=none,fillstyle=solid,fillcolor=curcolor]
{
\newpath
\moveto(68.39355469,289.07801687)
\lineto(70.77246094,289.07801687)
\lineto(70.77246094,281.61317312)
\lineto(68.39355469,281.61317312)
\closepath
\moveto(68.39355469,291.98192312)
\lineto(70.77246094,291.98192312)
\lineto(70.77246094,290.04598562)
\lineto(68.39355469,290.04598562)
\closepath
\moveto(68.39355469,291.98192312)
}
}
{
\newrgbcolor{curcolor}{0 0 0}
\pscustom[linestyle=none,fillstyle=solid,fillcolor=curcolor]
{
\newpath
\moveto(76.61308594,287.55223562)
\curveto(76.08808594,287.55223562)(75.68203099,287.35946192)(75.39902344,286.97801687)
\curveto(75.12421849,286.60477442)(74.98886719,286.05926687)(74.98886719,285.33739187)
\curveto(74.98886719,284.62372052)(75.12421849,284.07821192)(75.39902344,283.69676687)
\curveto(75.68203099,283.31122052)(76.08808594,283.12254812)(76.61308594,283.12254812)
\curveto(77.13808594,283.12254812)(77.53593724,283.31122052)(77.81074219,283.69676687)
\curveto(78.08144584,284.07821192)(78.22089844,284.62372052)(78.22089844,285.33739187)
\curveto(78.22089844,286.05926687)(78.08144584,286.60477442)(77.81074219,286.97801687)
\curveto(77.53593724,287.35946192)(77.13808594,287.55223562)(76.61308594,287.55223562)
\closepath
\moveto(76.61308594,289.25848562)
\curveto(77.90097604,289.25848562)(78.90996094,288.90575072)(79.63183594,288.20848562)
\curveto(80.35371094,287.51942312)(80.71464844,286.55965802)(80.71464844,285.33739187)
\curveto(80.71464844,284.11102442)(80.35371094,283.14715802)(79.63183594,282.44989187)
\curveto(78.90996094,281.76082937)(77.90097604,281.41629812)(76.61308594,281.41629812)
\curveto(75.33339844,281.41629812)(74.32441354,281.76082937)(73.59433594,282.44989187)
\curveto(72.87246094,283.14715802)(72.51152344,284.11102442)(72.51152344,285.33739187)
\curveto(72.51152344,286.55965802)(72.87246094,287.51942312)(73.59433594,288.20848562)
\curveto(74.32441354,288.90575072)(75.33339844,289.25848562)(76.61308594,289.25848562)
\closepath
\moveto(76.61308594,289.25848562)
}
}
{
\newrgbcolor{curcolor}{0 0 0}
\pscustom[linestyle=none,fillstyle=solid,fillcolor=curcolor]
{
\newpath
\moveto(89.9390625,286.15770437)
\lineto(89.9390625,281.61317312)
\lineto(87.54375,281.61317312)
\lineto(87.54375,285.09129812)
\curveto(87.54375,285.73524317)(87.52734375,286.17821192)(87.49453125,286.42020437)
\curveto(87.46992135,286.67039942)(87.42070365,286.85086817)(87.346875,286.96161062)
\curveto(87.25664115,287.11336817)(87.13359375,287.23231322)(86.96953125,287.32254812)
\curveto(86.81367135,287.40868067)(86.63320365,287.45379812)(86.428125,287.45379812)
\curveto(85.92363255,287.45379812)(85.52988255,287.26102442)(85.246875,286.87957937)
\curveto(84.95976615,286.49403302)(84.8203125,285.96493067)(84.8203125,285.28817312)
\lineto(84.8203125,281.61317312)
\lineto(82.44140625,281.61317312)
\lineto(82.44140625,289.07801687)
\lineto(84.8203125,289.07801687)
\lineto(84.8203125,287.97879812)
\curveto(85.18125,288.41356322)(85.56269505,288.73758692)(85.96875,288.94676687)
\curveto(86.37070365,289.15184552)(86.821875,289.25848562)(87.3140625,289.25848562)
\curveto(88.17539115,289.25848562)(88.8275388,288.98778302)(89.26640625,288.45457937)
\curveto(89.7134763,287.92957937)(89.9390625,287.16258692)(89.9390625,286.15770437)
\closepath
\moveto(89.9390625,286.15770437)
}
}
{
\newrgbcolor{curcolor}{0 0 0}
\pscustom[linestyle=none,fillstyle=solid,fillcolor=curcolor]
{
\newpath
\moveto(99.65976563,286.15770437)
\lineto(99.65976563,281.61317312)
\lineto(97.26445313,281.61317312)
\lineto(97.26445313,285.09129812)
\curveto(97.26445313,285.73524317)(97.24804688,286.17821192)(97.21523437,286.42020437)
\curveto(97.19062448,286.67039942)(97.14140678,286.85086817)(97.06757812,286.96161062)
\curveto(96.97734428,287.11336817)(96.85429688,287.23231322)(96.69023438,287.32254812)
\curveto(96.53437448,287.40868067)(96.35390677,287.45379812)(96.14882812,287.45379812)
\curveto(95.64433567,287.45379812)(95.25058567,287.26102442)(94.96757813,286.87957937)
\curveto(94.68046928,286.49403302)(94.54101562,285.96493067)(94.54101562,285.28817312)
\lineto(94.54101562,281.61317312)
\lineto(92.16210938,281.61317312)
\lineto(92.16210938,289.07801687)
\lineto(94.54101562,289.07801687)
\lineto(94.54101562,287.97879812)
\curveto(94.90195313,288.41356322)(95.28339818,288.73758692)(95.68945312,288.94676687)
\curveto(96.09140677,289.15184552)(96.54257813,289.25848562)(97.03476563,289.25848562)
\curveto(97.89609427,289.25848562)(98.54824192,288.98778302)(98.98710938,288.45457937)
\curveto(99.43417942,287.92957937)(99.65976563,287.16258692)(99.65976563,286.15770437)
\closepath
\moveto(99.65976563,286.15770437)
}
}
{
\newrgbcolor{curcolor}{0 0 0}
\pscustom[linestyle=none,fillstyle=solid,fillcolor=curcolor]
{
\newpath
\moveto(105.2296875,284.97645437)
\curveto(104.72519505,284.97645437)(104.3478513,284.88622052)(104.09765625,284.71395437)
\curveto(103.8433599,284.54989187)(103.7203125,284.30379812)(103.7203125,283.97567312)
\curveto(103.7203125,283.66805567)(103.81875,283.42606322)(104.015625,283.25379812)
\curveto(104.2207026,283.08973562)(104.5078125,283.00770437)(104.86875,283.00770437)
\curveto(105.31582005,283.00770437)(105.6931638,283.16356322)(106.00078125,283.48348562)
\curveto(106.3042974,283.81161062)(106.46015625,284.21356322)(106.46015625,284.69754812)
\lineto(106.46015625,284.97645437)
\closepath
\moveto(108.871875,285.87879812)
\lineto(108.871875,281.61317312)
\lineto(106.46015625,281.61317312)
\lineto(106.46015625,282.71239187)
\curveto(106.1402349,282.26122052)(105.7792974,281.93309552)(105.37734375,281.72801687)
\curveto(104.98359375,281.52293875)(104.4996099,281.41629812)(103.93359375,281.41629812)
\curveto(103.1542974,281.41629812)(102.52675755,281.6377825)(102.046875,282.08895437)
\curveto(101.5628901,282.54832937)(101.325,283.13895437)(101.325,283.86082937)
\curveto(101.325,284.73446192)(101.6244138,285.37430567)(102.22734375,285.78036062)
\curveto(102.8261724,286.18231322)(103.77363255,286.38739187)(105.065625,286.38739187)
\lineto(106.46015625,286.38739187)
\lineto(106.46015625,286.58426687)
\curveto(106.46015625,286.95340802)(106.3125,287.22821192)(106.0171875,287.40457937)
\curveto(105.721875,287.57684552)(105.2542974,287.66707937)(104.62265625,287.66707937)
\curveto(104.1181638,287.66707937)(103.64238255,287.61786062)(103.1953125,287.51942312)
\curveto(102.75644505,287.42098562)(102.3503901,287.26512572)(101.98125,287.06004812)
\lineto(101.98125,288.88114187)
\curveto(102.4816401,289.00008692)(102.9902349,289.09442312)(103.50703125,289.16004812)
\curveto(104.03203125,289.22567312)(104.5488276,289.25848562)(105.065625,289.25848562)
\curveto(106.4109375,289.25848562)(107.37890625,288.98778302)(107.96953125,288.45457937)
\curveto(108.5683599,287.92957937)(108.871875,287.06825072)(108.871875,285.87879812)
\closepath
\moveto(108.871875,285.87879812)
}
}
{
\newrgbcolor{curcolor}{0 0 0}
\pscustom[linestyle=none,fillstyle=solid,fillcolor=curcolor]
{
\newpath
\moveto(111.09082031,289.07801687)
\lineto(113.46972656,289.07801687)
\lineto(113.46972656,281.61317312)
\lineto(111.09082031,281.61317312)
\closepath
\moveto(111.09082031,291.98192312)
\lineto(113.46972656,291.98192312)
\lineto(113.46972656,290.04598562)
\lineto(111.09082031,290.04598562)
\closepath
\moveto(111.09082031,291.98192312)
}
}
{
\newrgbcolor{curcolor}{0 0 0}
\pscustom[linestyle=none,fillstyle=solid,fillcolor=curcolor]
{
\newpath
\moveto(121.31191406,287.04364187)
\curveto(121.10273411,287.14207937)(120.89355416,287.21180567)(120.68847656,287.25692312)
\curveto(120.47929661,287.31024317)(120.27011666,287.33895437)(120.06503906,287.33895437)
\curveto(119.44980416,287.33895437)(118.97402396,287.14207937)(118.63769531,286.74832937)
\curveto(118.30957031,286.35457937)(118.14550781,285.78856322)(118.14550781,285.05848562)
\lineto(118.14550781,281.61317312)
\lineto(115.76660156,281.61317312)
\lineto(115.76660156,289.07801687)
\lineto(118.14550781,289.07801687)
\lineto(118.14550781,287.84754812)
\curveto(118.44902396,288.33973562)(118.80175781,288.69247052)(119.19550781,288.91395437)
\curveto(119.59746146,289.14364187)(120.08144531,289.25848562)(120.63925781,289.25848562)
\curveto(120.71308646,289.25848562)(120.79511666,289.25028302)(120.88535156,289.24207937)
\curveto(120.98378906,289.24207937)(121.11914036,289.22977442)(121.29550781,289.20926687)
\closepath
\moveto(121.31191406,287.04364187)
}
}
{
\newrgbcolor{curcolor}{0 0 0}
\pscustom[linestyle=none,fillstyle=solid,fillcolor=curcolor]
{
\newpath
\moveto(129.94160156,285.37020437)
\lineto(129.94160156,284.68114187)
\lineto(124.36347656,284.68114187)
\curveto(124.41679661,284.12332937)(124.61777396,283.70086817)(124.97050781,283.41786062)
\curveto(125.31914036,283.14305567)(125.80722656,283.00770437)(126.43066406,283.00770437)
\curveto(126.92285156,283.00770437)(127.43144531,283.07743067)(127.95644531,283.22098562)
\curveto(128.48964896,283.37274317)(129.03925781,283.60243067)(129.59707031,283.91004812)
\lineto(129.59707031,282.07254812)
\curveto(129.02695286,281.85106322)(128.45683646,281.68700125)(127.89082031,281.58036062)
\curveto(127.33300781,281.47372052)(126.76699166,281.41629812)(126.20097656,281.41629812)
\curveto(124.85566406,281.41629812)(123.80566406,281.76082937)(123.05097656,282.44989187)
\curveto(122.30449166,283.13895437)(121.93535156,284.09872052)(121.93535156,285.33739187)
\curveto(121.93535156,286.55145437)(122.30039036,287.50711817)(123.03457031,288.20848562)
\curveto(123.76464896,288.90575072)(124.77773411,289.25848562)(126.06972656,289.25848562)
\curveto(127.23867161,289.25848562)(128.17382786,288.90164942)(128.87519531,288.19207937)
\curveto(129.58476536,287.49071192)(129.94160156,286.55145437)(129.94160156,285.37020437)
\closepath
\moveto(127.49707031,286.15770437)
\curveto(127.49707031,286.61707937)(127.35761666,286.98211817)(127.08691406,287.25692312)
\curveto(126.82441406,287.52762572)(126.47988281,287.66707937)(126.05332031,287.66707937)
\curveto(125.59394531,287.66707937)(125.21660156,287.53582937)(124.92128906,287.27332937)
\curveto(124.63417916,287.01903302)(124.46191406,286.64989187)(124.39628906,286.15770437)
\closepath
\moveto(127.49707031,286.15770437)
}
}
{
\newrgbcolor{curcolor}{0 0 0}
\pscustom[linewidth=2.10000003,linecolor=curcolor]
{
\newpath
\moveto(444.096681,305.10692575)
\lineto(594.615819,305.10692575)
\lineto(594.615819,265.20692575)
\lineto(444.096681,265.20692575)
\closepath
\moveto(444.096681,305.10692575)
}
}
{
\newrgbcolor{curcolor}{0 0 0}
\pscustom[linestyle=none,fillstyle=solid,fillcolor=curcolor]
{
\newpath
\moveto(487.2,286.60887625)
\curveto(487.73320365,286.60887625)(488.11875,286.70731375)(488.3484375,286.90418875)
\curveto(488.578125,287.10106375)(488.69296875,287.42918875)(488.69296875,287.88856375)
\curveto(488.69296875,288.3356338)(488.578125,288.6596575)(488.3484375,288.8565325)
\curveto(488.11875,289.0534075)(487.73320365,289.151845)(487.2,289.151845)
\lineto(486.1171875,289.151845)
\lineto(486.1171875,286.60887625)
\closepath
\moveto(486.1171875,284.83700125)
\lineto(486.1171875,281.06356375)
\lineto(483.54140625,281.06356375)
\lineto(483.54140625,291.0221575)
\lineto(487.4625,291.0221575)
\curveto(488.775,291.0221575)(489.73476615,290.7965713)(490.35,290.34950125)
\curveto(490.96113255,289.9106338)(491.26875,289.21747)(491.26875,288.2659075)
\curveto(491.26875,287.6096575)(491.1087888,287.06825125)(490.79296875,286.64168875)
\curveto(490.47304635,286.22332885)(489.9931638,285.91981375)(489.34921875,285.72293875)
\curveto(489.6978513,285.63270385)(490.00957005,285.4481338)(490.284375,285.16512625)
\curveto(490.56738255,284.8903213)(490.85039115,284.4637588)(491.1375,283.88543875)
\lineto(492.53203125,281.06356375)
\lineto(489.7921875,281.06356375)
\lineto(488.59453125,283.5409075)
\curveto(488.34023385,284.033095)(488.0900388,284.36942365)(487.83984375,284.558095)
\curveto(487.58554635,284.74266505)(487.25332005,284.83700125)(486.8390625,284.83700125)
\closepath
\moveto(486.1171875,284.83700125)
}
}
{
\newrgbcolor{curcolor}{0 0 0}
\pscustom[linestyle=none,fillstyle=solid,fillcolor=curcolor]
{
\newpath
\moveto(501.41191406,284.820595)
\lineto(501.41191406,284.1315325)
\lineto(495.83378906,284.1315325)
\curveto(495.88710911,283.57372)(496.08808646,283.1512588)(496.44082031,282.86825125)
\curveto(496.78945286,282.5934463)(497.27753906,282.458095)(497.90097656,282.458095)
\curveto(498.39316406,282.458095)(498.90175781,282.5278213)(499.42675781,282.67137625)
\curveto(499.95996146,282.8231338)(500.50957031,283.0528213)(501.06738281,283.36043875)
\lineto(501.06738281,281.52293875)
\curveto(500.49726536,281.30145385)(499.92714896,281.13739187)(499.36113281,281.03075125)
\curveto(498.80332031,280.92411115)(498.23730416,280.86668875)(497.67128906,280.86668875)
\curveto(496.32597656,280.86668875)(495.27597656,281.21122)(494.52128906,281.9002825)
\curveto(493.77480416,282.589345)(493.40566406,283.54911115)(493.40566406,284.7877825)
\curveto(493.40566406,286.001845)(493.77070286,286.9575088)(494.50488281,287.65887625)
\curveto(495.23496146,288.35614135)(496.24804661,288.70887625)(497.54003906,288.70887625)
\curveto(498.70898411,288.70887625)(499.64414036,288.35204005)(500.34550781,287.64247)
\curveto(501.05507786,286.94110255)(501.41191406,286.001845)(501.41191406,284.820595)
\closepath
\moveto(498.96738281,285.608095)
\curveto(498.96738281,286.06747)(498.82792916,286.4325088)(498.55722656,286.70731375)
\curveto(498.29472656,286.97801635)(497.95019531,287.11747)(497.52363281,287.11747)
\curveto(497.06425781,287.11747)(496.68691406,286.98622)(496.39160156,286.72372)
\curveto(496.10449166,286.46942365)(495.93222656,286.1002825)(495.86660156,285.608095)
\closepath
\moveto(498.96738281,285.608095)
}
}
{
\newrgbcolor{curcolor}{0 0 0}
\pscustom[linestyle=none,fillstyle=solid,fillcolor=curcolor]
{
\newpath
\moveto(510.1359375,287.2815325)
\curveto(510.43125,287.74911115)(510.7839849,288.1059463)(511.20234375,288.34793875)
\curveto(511.6166013,288.58582885)(512.0759763,288.70887625)(512.58046875,288.70887625)
\curveto(513.43359375,288.70887625)(514.0816401,288.43817365)(514.5328125,287.90497)
\curveto(514.97988255,287.37997)(515.20546875,286.61297755)(515.20546875,285.608095)
\lineto(515.20546875,281.06356375)
\lineto(512.81015625,281.06356375)
\lineto(512.81015625,284.951845)
\lineto(512.81015625,285.13231375)
\curveto(512.8183599,285.19793875)(512.8265625,285.2840713)(512.8265625,285.39481375)
\curveto(512.8265625,285.92801635)(512.74453125,286.31356375)(512.58046875,286.54325125)
\curveto(512.4246099,286.78114135)(512.1785151,286.90418875)(511.8421875,286.90418875)
\curveto(511.3828125,286.90418875)(511.0259763,286.71551635)(510.77578125,286.34637625)
\curveto(510.5337888,285.9731338)(510.4066401,285.43172755)(510.3984375,284.7221575)
\lineto(510.3984375,281.06356375)
\lineto(508.003125,281.06356375)
\lineto(508.003125,284.951845)
\curveto(508.003125,285.78036115)(507.9292974,286.31356375)(507.78984375,286.54325125)
\curveto(507.6462888,286.78114135)(507.39609375,286.90418875)(507.03515625,286.90418875)
\curveto(506.57578125,286.90418875)(506.21894505,286.71551635)(505.96875,286.34637625)
\curveto(505.7144526,285.9731338)(505.59140625,285.43582885)(505.59140625,284.73856375)
\lineto(505.59140625,281.06356375)
\lineto(503.19609375,281.06356375)
\lineto(503.19609375,288.5284075)
\lineto(505.59140625,288.5284075)
\lineto(505.59140625,287.42918875)
\curveto(505.88671875,287.85575125)(506.2230474,288.1715713)(506.60859375,288.38075125)
\curveto(506.9900388,288.5981338)(507.4042974,288.70887625)(507.85546875,288.70887625)
\curveto(508.38046875,288.70887625)(508.83984375,288.58172755)(509.23359375,288.3315325)
\curveto(509.62734375,288.07723615)(509.92675755,287.72860255)(510.1359375,287.2815325)
\closepath
\moveto(510.1359375,287.2815325)
}
}
{
\newrgbcolor{curcolor}{0 0 0}
\pscustom[linestyle=none,fillstyle=solid,fillcolor=curcolor]
{
\newpath
\moveto(517.44492187,288.5284075)
\lineto(519.82382812,288.5284075)
\lineto(519.82382812,281.06356375)
\lineto(517.44492187,281.06356375)
\closepath
\moveto(517.44492187,291.43231375)
\lineto(519.82382812,291.43231375)
\lineto(519.82382812,289.49637625)
\lineto(517.44492187,289.49637625)
\closepath
\moveto(517.44492187,291.43231375)
}
}
{
\newrgbcolor{curcolor}{0 0 0}
\pscustom[linestyle=none,fillstyle=solid,fillcolor=curcolor]
{
\newpath
\moveto(529.61835937,285.608095)
\lineto(529.61835937,281.06356375)
\lineto(527.22304688,281.06356375)
\lineto(527.22304688,284.54168875)
\curveto(527.22304688,285.1856338)(527.20664062,285.62860255)(527.17382812,285.870595)
\curveto(527.14921823,286.12079005)(527.10000052,286.3012588)(527.02617187,286.41200125)
\curveto(526.93593802,286.5637588)(526.81289063,286.68270385)(526.64882813,286.77293875)
\curveto(526.49296823,286.8590713)(526.31250053,286.90418875)(526.10742188,286.90418875)
\curveto(525.60292942,286.90418875)(525.20917943,286.71141505)(524.92617188,286.32997)
\curveto(524.63906302,285.94442365)(524.49960937,285.4153213)(524.49960937,284.73856375)
\lineto(524.49960937,281.06356375)
\lineto(522.12070312,281.06356375)
\lineto(522.12070312,288.5284075)
\lineto(524.49960937,288.5284075)
\lineto(524.49960937,287.42918875)
\curveto(524.86054687,287.86395385)(525.24199192,288.18797755)(525.64804687,288.3971575)
\curveto(526.05000052,288.60223615)(526.50117187,288.70887625)(526.99335937,288.70887625)
\curveto(527.85468802,288.70887625)(528.50683567,288.43817365)(528.94570313,287.90497)
\curveto(529.39277318,287.37997)(529.61835937,286.61297755)(529.61835937,285.608095)
\closepath
\moveto(529.61835937,285.608095)
}
}
{
\newrgbcolor{curcolor}{0 0 0}
\pscustom[linestyle=none,fillstyle=solid,fillcolor=curcolor]
{
\newpath
\moveto(536.9109375,287.42918875)
\lineto(536.9109375,291.43231375)
\lineto(539.32265625,291.43231375)
\lineto(539.32265625,281.06356375)
\lineto(536.9109375,281.06356375)
\lineto(536.9109375,282.14637625)
\curveto(536.5828125,281.7075088)(536.221875,281.38348615)(535.828125,281.1784075)
\curveto(535.434375,280.97332937)(534.9791013,280.86668875)(534.46640625,280.86668875)
\curveto(533.5353513,280.86668875)(532.7765625,281.23172755)(532.1859375,281.9659075)
\curveto(531.5953125,282.69598615)(531.3,283.639345)(531.3,284.7877825)
\curveto(531.3,285.93622)(531.5953125,286.87547755)(532.1859375,287.6096575)
\curveto(532.7765625,288.33973615)(533.5353513,288.70887625)(534.46640625,288.70887625)
\curveto(534.9791013,288.70887625)(535.434375,288.60223615)(535.828125,288.3971575)
\curveto(536.221875,288.18797755)(536.5828125,287.86395385)(536.9109375,287.42918875)
\closepath
\moveto(535.35234375,282.60575125)
\curveto(535.85273385,282.60575125)(536.23828125,282.7903213)(536.50078125,283.16356375)
\curveto(536.77148385,283.53270385)(536.9109375,284.07411115)(536.9109375,284.7877825)
\curveto(536.9109375,285.49735255)(536.77148385,286.0387588)(536.50078125,286.41200125)
\curveto(536.23828125,286.78114135)(535.85273385,286.96981375)(535.35234375,286.96981375)
\curveto(534.83554635,286.96981375)(534.44179635,286.78114135)(534.17109375,286.41200125)
\curveto(533.90859375,286.0387588)(533.77734375,285.49735255)(533.77734375,284.7877825)
\curveto(533.77734375,284.07411115)(533.90859375,283.53270385)(534.17109375,283.16356375)
\curveto(534.44179635,282.7903213)(534.83554635,282.60575125)(535.35234375,282.60575125)
\closepath
\moveto(535.35234375,282.60575125)
}
}
{
\newrgbcolor{curcolor}{0 0 0}
\pscustom[linestyle=none,fillstyle=solid,fillcolor=curcolor]
{
\newpath
\moveto(549.0515625,284.820595)
\lineto(549.0515625,284.1315325)
\lineto(543.4734375,284.1315325)
\curveto(543.52675755,283.57372)(543.7277349,283.1512588)(544.08046875,282.86825125)
\curveto(544.4291013,282.5934463)(544.9171875,282.458095)(545.540625,282.458095)
\curveto(546.0328125,282.458095)(546.54140625,282.5278213)(547.06640625,282.67137625)
\curveto(547.5996099,282.8231338)(548.14921875,283.0528213)(548.70703125,283.36043875)
\lineto(548.70703125,281.52293875)
\curveto(548.1369138,281.30145385)(547.5667974,281.13739187)(547.00078125,281.03075125)
\curveto(546.44296875,280.92411115)(545.8769526,280.86668875)(545.3109375,280.86668875)
\curveto(543.965625,280.86668875)(542.915625,281.21122)(542.1609375,281.9002825)
\curveto(541.4144526,282.589345)(541.0453125,283.54911115)(541.0453125,284.7877825)
\curveto(541.0453125,286.001845)(541.4103513,286.9575088)(542.14453125,287.65887625)
\curveto(542.8746099,288.35614135)(543.88769505,288.70887625)(545.1796875,288.70887625)
\curveto(546.34863255,288.70887625)(547.2837888,288.35204005)(547.98515625,287.64247)
\curveto(548.6947263,286.94110255)(549.0515625,286.001845)(549.0515625,284.820595)
\closepath
\moveto(546.60703125,285.608095)
\curveto(546.60703125,286.06747)(546.4675776,286.4325088)(546.196875,286.70731375)
\curveto(545.934375,286.97801635)(545.58984375,287.11747)(545.16328125,287.11747)
\curveto(544.70390625,287.11747)(544.3265625,286.98622)(544.03125,286.72372)
\curveto(543.7441401,286.46942365)(543.571875,286.1002825)(543.50625,285.608095)
\closepath
\moveto(546.60703125,285.608095)
}
}
{
\newrgbcolor{curcolor}{0 0 0}
\pscustom[linestyle=none,fillstyle=solid,fillcolor=curcolor]
{
\newpath
\moveto(556.39746094,286.4940325)
\curveto(556.18828099,286.59247)(555.97910104,286.6621963)(555.77402344,286.70731375)
\curveto(555.56484349,286.7606338)(555.35566354,286.789345)(555.15058594,286.789345)
\curveto(554.53535104,286.789345)(554.05957084,286.59247)(553.72324219,286.19872)
\curveto(553.39511719,285.80497)(553.23105469,285.23895385)(553.23105469,284.50887625)
\lineto(553.23105469,281.06356375)
\lineto(550.85214844,281.06356375)
\lineto(550.85214844,288.5284075)
\lineto(553.23105469,288.5284075)
\lineto(553.23105469,287.29793875)
\curveto(553.53457084,287.79012625)(553.88730469,288.14286115)(554.28105469,288.364345)
\curveto(554.68300834,288.5940325)(555.16699219,288.70887625)(555.72480469,288.70887625)
\curveto(555.79863334,288.70887625)(555.88066354,288.70067365)(555.97089844,288.69247)
\curveto(556.06933594,288.69247)(556.20468724,288.68016505)(556.38105469,288.6596575)
\closepath
\moveto(556.39746094,286.4940325)
}
}
{
\newrgbcolor{curcolor}{0 0 0}
\pscustom[linewidth=2.10000003,linecolor=curcolor]
{
\newpath
\moveto(142.33653,135.99539575)
\curveto(142.33653,117.66552475)(112.616595,102.80554675)(75.956832,102.80554675)
\curveto(39.297069,102.80554675)(9.577155,117.66552475)(9.577155,135.99539575)
\curveto(9.577155,154.32528775)(39.297069,169.18524475)(75.956832,169.18524475)
\curveto(112.616595,169.18524475)(142.33653,154.32528775)(142.33653,135.99539575)
}
}
{
\newrgbcolor{curcolor}{0 0 0}
\pscustom[linestyle=none,fillstyle=solid,fillcolor=curcolor]
{
\newpath
\moveto(50.6953125,133.7190325)
\lineto(46.6921875,133.7190325)
\lineto(46.05234375,131.89793875)
\lineto(43.4765625,131.89793875)
\lineto(47.16796875,141.8565325)
\lineto(50.21953125,141.8565325)
\lineto(53.9109375,131.89793875)
\lineto(51.33515625,131.89793875)
\closepath
\moveto(47.33203125,135.5565325)
\lineto(50.05546875,135.5565325)
\lineto(48.69375,139.51043875)
\closepath
\moveto(47.33203125,135.5565325)
}
}
{
\newrgbcolor{curcolor}{0 0 0}
\pscustom[linestyle=none,fillstyle=solid,fillcolor=curcolor]
{
\newpath
\moveto(62.61855469,136.44247)
\lineto(62.61855469,131.89793875)
\lineto(60.22324219,131.89793875)
\lineto(60.22324219,135.37606375)
\curveto(60.22324219,136.0200088)(60.20683594,136.46297755)(60.17402344,136.70497)
\curveto(60.14941354,136.95516505)(60.10019584,137.1356338)(60.02636719,137.24637625)
\curveto(59.93613334,137.3981338)(59.81308594,137.51707885)(59.64902344,137.60731375)
\curveto(59.49316354,137.6934463)(59.31269584,137.73856375)(59.10761719,137.73856375)
\curveto(58.60312474,137.73856375)(58.20937474,137.54579005)(57.92636719,137.164345)
\curveto(57.63925834,136.77879865)(57.49980469,136.2496963)(57.49980469,135.57293875)
\lineto(57.49980469,131.89793875)
\lineto(55.12089844,131.89793875)
\lineto(55.12089844,139.3627825)
\lineto(57.49980469,139.3627825)
\lineto(57.49980469,138.26356375)
\curveto(57.86074219,138.69832885)(58.24218724,139.02235255)(58.64824219,139.2315325)
\curveto(59.05019584,139.43661115)(59.50136719,139.54325125)(59.99355469,139.54325125)
\curveto(60.85488334,139.54325125)(61.50703099,139.27254865)(61.94589844,138.739345)
\curveto(62.39296849,138.214345)(62.61855469,137.44735255)(62.61855469,136.44247)
\closepath
\moveto(62.61855469,136.44247)
}
}
{
\newrgbcolor{curcolor}{0 0 0}
\pscustom[linestyle=none,fillstyle=solid,fillcolor=curcolor]
{
\newpath
\moveto(70.66582031,139.133095)
\lineto(70.66582031,137.31200125)
\curveto(70.16132786,137.5293838)(69.66914036,137.689345)(69.18925781,137.7877825)
\curveto(68.71757786,137.89442365)(68.27050781,137.951845)(67.84394531,137.951845)
\curveto(67.39277396,137.951845)(67.05644531,137.89442365)(66.82675781,137.7877825)
\curveto(66.60527396,137.67704005)(66.49863281,137.50067365)(66.49863281,137.2627825)
\curveto(66.49863281,137.07411115)(66.58066406,136.92645385)(66.74472656,136.81981375)
\curveto(66.91699166,136.72137625)(67.22050781,136.6434463)(67.64707031,136.59012625)
\lineto(68.07363281,136.5409075)
\curveto(69.29589896,136.38504865)(70.11621146,136.13075125)(70.53457031,135.76981375)
\curveto(70.96113281,135.41707885)(71.17441406,134.85926635)(71.17441406,134.09637625)
\curveto(71.17441406,133.2965713)(70.87910156,132.69364135)(70.28847656,132.29168875)
\curveto(69.70605416,131.89793875)(68.83242161,131.70106375)(67.66347656,131.70106375)
\curveto(67.17128906,131.70106375)(66.66269531,131.74207885)(66.13769531,131.8159075)
\curveto(65.61269531,131.89383719)(65.07128906,132.0127825)(64.51347656,132.176845)
\lineto(64.51347656,133.9815325)
\curveto(64.98105416,133.751845)(65.46914036,133.57547755)(65.97363281,133.4565325)
\curveto(66.47402396,133.34579005)(66.98261666,133.29247)(67.49941406,133.29247)
\curveto(67.96699166,133.29247)(68.31972656,133.34989135)(68.54941406,133.47293875)
\curveto(68.78730416,133.60418875)(68.91035156,133.80106375)(68.91035156,134.06356375)
\curveto(68.91035156,134.26864135)(68.82832031,134.42450125)(68.66425781,134.52293875)
\curveto(68.50019531,134.62957885)(68.17617161,134.71981375)(67.69628906,134.78543875)
\lineto(67.26972656,134.8346575)
\curveto(66.20742161,134.9659075)(65.46503906,135.21200125)(65.03847656,135.57293875)
\curveto(64.61191406,135.93387625)(64.39863281,136.4793838)(64.39863281,137.21356375)
\curveto(64.39863281,138.00106375)(64.66523411,138.58348615)(65.20253906,138.9690325)
\curveto(65.74804661,139.35047755)(66.58476536,139.54325125)(67.71269531,139.54325125)
\curveto(68.14746146,139.54325125)(68.60683646,139.51043875)(69.09082031,139.44481375)
\curveto(69.57070286,139.37918875)(70.09570286,139.27254865)(70.66582031,139.133095)
\closepath
\moveto(70.66582031,139.133095)
}
}
{
\newrgbcolor{curcolor}{0 0 0}
\pscustom[linestyle=none,fillstyle=solid,fillcolor=curcolor]
{
\newpath
\moveto(72.29003906,139.3627825)
\lineto(74.61972656,139.3627825)
\lineto(75.86660156,134.21122)
\lineto(77.12988281,139.3627825)
\lineto(79.11503906,139.3627825)
\lineto(80.37832031,134.276845)
\lineto(81.62519531,139.3627825)
\lineto(83.95488281,139.3627825)
\lineto(81.98613281,131.89793875)
\lineto(79.37753906,131.89793875)
\lineto(78.11425781,137.033095)
\lineto(76.86738281,131.89793875)
\lineto(74.25878906,131.89793875)
\closepath
\moveto(72.29003906,139.3627825)
}
}
{
\newrgbcolor{curcolor}{0 0 0}
\pscustom[linestyle=none,fillstyle=solid,fillcolor=curcolor]
{
\newpath
\moveto(93.0234375,135.65497)
\lineto(93.0234375,134.9659075)
\lineto(87.4453125,134.9659075)
\curveto(87.49863255,134.408095)(87.6996099,133.9856338)(88.05234375,133.70262625)
\curveto(88.4009763,133.4278213)(88.8890625,133.29247)(89.5125,133.29247)
\curveto(90.0046875,133.29247)(90.51328125,133.3621963)(91.03828125,133.50575125)
\curveto(91.5714849,133.6575088)(92.12109375,133.8871963)(92.67890625,134.19481375)
\lineto(92.67890625,132.35731375)
\curveto(92.1087888,132.13582885)(91.5386724,131.97176687)(90.97265625,131.86512625)
\curveto(90.41484375,131.75848615)(89.8488276,131.70106375)(89.2828125,131.70106375)
\curveto(87.9375,131.70106375)(86.8875,132.045595)(86.1328125,132.7346575)
\curveto(85.3863276,133.42372)(85.0171875,134.38348615)(85.0171875,135.6221575)
\curveto(85.0171875,136.83622)(85.3822263,137.7918838)(86.11640625,138.49325125)
\curveto(86.8464849,139.19051635)(87.85957005,139.54325125)(89.1515625,139.54325125)
\curveto(90.32050755,139.54325125)(91.2556638,139.18641505)(91.95703125,138.476845)
\curveto(92.6666013,137.77547755)(93.0234375,136.83622)(93.0234375,135.65497)
\closepath
\moveto(90.57890625,136.44247)
\curveto(90.57890625,136.901845)(90.4394526,137.2668838)(90.16875,137.54168875)
\curveto(89.90625,137.81239135)(89.56171875,137.951845)(89.13515625,137.951845)
\curveto(88.67578125,137.951845)(88.2984375,137.820595)(88.003125,137.558095)
\curveto(87.7160151,137.30379865)(87.54375,136.9346575)(87.478125,136.44247)
\closepath
\moveto(90.57890625,136.44247)
}
}
{
\newrgbcolor{curcolor}{0 0 0}
\pscustom[linestyle=none,fillstyle=solid,fillcolor=curcolor]
{
\newpath
\moveto(100.36933594,137.3284075)
\curveto(100.16015599,137.426845)(99.95097604,137.4965713)(99.74589844,137.54168875)
\curveto(99.53671849,137.5950088)(99.32753854,137.62372)(99.12246094,137.62372)
\curveto(98.50722604,137.62372)(98.03144584,137.426845)(97.69511719,137.033095)
\curveto(97.36699219,136.639345)(97.20292969,136.07332885)(97.20292969,135.34325125)
\lineto(97.20292969,131.89793875)
\lineto(94.82402344,131.89793875)
\lineto(94.82402344,139.3627825)
\lineto(97.20292969,139.3627825)
\lineto(97.20292969,138.13231375)
\curveto(97.50644584,138.62450125)(97.85917969,138.97723615)(98.25292969,139.19872)
\curveto(98.65488334,139.4284075)(99.13886719,139.54325125)(99.69667969,139.54325125)
\curveto(99.77050834,139.54325125)(99.85253854,139.53504865)(99.94277344,139.526845)
\curveto(100.04121094,139.526845)(100.17656224,139.51454005)(100.35292969,139.4940325)
\closepath
\moveto(100.36933594,137.3284075)
}
}
{
\newrgbcolor{curcolor}{0 0 0}
\pscustom[linestyle=none,fillstyle=solid,fillcolor=curcolor]
{
\newpath
\moveto(107.37480469,139.133095)
\lineto(107.37480469,137.31200125)
\curveto(106.87031224,137.5293838)(106.37812474,137.689345)(105.89824219,137.7877825)
\curveto(105.42656224,137.89442365)(104.97949219,137.951845)(104.55292969,137.951845)
\curveto(104.10175834,137.951845)(103.76542969,137.89442365)(103.53574219,137.7877825)
\curveto(103.31425834,137.67704005)(103.20761719,137.50067365)(103.20761719,137.2627825)
\curveto(103.20761719,137.07411115)(103.28964844,136.92645385)(103.45371094,136.81981375)
\curveto(103.62597604,136.72137625)(103.92949219,136.6434463)(104.35605469,136.59012625)
\lineto(104.78261719,136.5409075)
\curveto(106.00488334,136.38504865)(106.82519584,136.13075125)(107.24355469,135.76981375)
\curveto(107.67011719,135.41707885)(107.88339844,134.85926635)(107.88339844,134.09637625)
\curveto(107.88339844,133.2965713)(107.58808594,132.69364135)(106.99746094,132.29168875)
\curveto(106.41503854,131.89793875)(105.54140599,131.70106375)(104.37246094,131.70106375)
\curveto(103.88027344,131.70106375)(103.37167969,131.74207885)(102.84667969,131.8159075)
\curveto(102.32167969,131.89383719)(101.78027344,132.0127825)(101.22246094,132.176845)
\lineto(101.22246094,133.9815325)
\curveto(101.69003854,133.751845)(102.17812474,133.57547755)(102.68261719,133.4565325)
\curveto(103.18300834,133.34579005)(103.69160104,133.29247)(104.20839844,133.29247)
\curveto(104.67597604,133.29247)(105.02871094,133.34989135)(105.25839844,133.47293875)
\curveto(105.49628854,133.60418875)(105.61933594,133.80106375)(105.61933594,134.06356375)
\curveto(105.61933594,134.26864135)(105.53730469,134.42450125)(105.37324219,134.52293875)
\curveto(105.20917969,134.62957885)(104.88515599,134.71981375)(104.40527344,134.78543875)
\lineto(103.97871094,134.8346575)
\curveto(102.91640599,134.9659075)(102.17402344,135.21200125)(101.74746094,135.57293875)
\curveto(101.32089844,135.93387625)(101.10761719,136.4793838)(101.10761719,137.21356375)
\curveto(101.10761719,138.00106375)(101.37421849,138.58348615)(101.91152344,138.9690325)
\curveto(102.45703099,139.35047755)(103.29374974,139.54325125)(104.42167969,139.54325125)
\curveto(104.85644584,139.54325125)(105.31582084,139.51043875)(105.79980469,139.44481375)
\curveto(106.27968724,139.37918875)(106.80468724,139.27254865)(107.37480469,139.133095)
\closepath
\moveto(107.37480469,139.133095)
}
}
{
\newrgbcolor{curcolor}{0 0 0}
\pscustom[linewidth=2.10000003,linecolor=curcolor]
{
\newpath
\moveto(401.830086,130.66746475)
\curveto(401.830086,112.33759375)(372.110151,97.47763675)(335.450388,97.47763675)
\curveto(298.790625,97.47763675)(269.070711,112.33759375)(269.070711,130.66746475)
\curveto(269.070711,148.99735675)(298.790625,163.85731375)(335.450388,163.85731375)
\curveto(372.110151,163.85731375)(401.830086,148.99735675)(401.830086,130.66746475)
}
}
{
\newrgbcolor{curcolor}{0 0 0}
\pscustom[linestyle=none,fillstyle=solid,fillcolor=curcolor]
{
\newpath
\moveto(298.45839844,128.39520437)
\lineto(294.45527344,128.39520437)
\lineto(293.81542969,126.57411062)
\lineto(291.23964844,126.57411062)
\lineto(294.93105469,136.53270437)
\lineto(297.98261719,136.53270437)
\lineto(301.67402344,126.57411062)
\lineto(299.09824219,126.57411062)
\closepath
\moveto(295.09511719,130.23270437)
\lineto(297.81855469,130.23270437)
\lineto(296.45683594,134.18661062)
\closepath
\moveto(295.09511719,130.23270437)
}
}
{
\newrgbcolor{curcolor}{0 0 0}
\pscustom[linestyle=none,fillstyle=solid,fillcolor=curcolor]
{
\newpath
\moveto(308.92148438,133.80926687)
\lineto(308.92148438,131.85692312)
\curveto(308.59335938,132.07430567)(308.26523438,132.23836817)(307.93710938,132.34911062)
\curveto(307.60898438,132.45575072)(307.26855442,132.51317312)(306.91992188,132.51317312)
\curveto(306.25136693,132.51317312)(305.73046928,132.31629812)(305.36132812,131.92254812)
\curveto(304.98808568,131.53700072)(304.80351562,130.99559552)(304.80351562,130.29832937)
\curveto(304.80351562,129.59696192)(304.98808568,129.05145437)(305.36132812,128.65770437)
\curveto(305.73046928,128.27215802)(306.25136693,128.08348562)(306.91992188,128.08348562)
\curveto(307.28906198,128.08348562)(307.64589818,128.13680567)(307.98632812,128.24754812)
\curveto(308.32265677,128.36649317)(308.63437448,128.53055567)(308.92148438,128.73973562)
\lineto(308.92148438,126.78739187)
\curveto(308.53593698,126.65614187)(308.15449193,126.55770437)(307.77304688,126.49207937)
\curveto(307.38749948,126.41825072)(307.00605443,126.37723562)(306.62460938,126.37723562)
\curveto(305.27929688,126.37723562)(304.22109427,126.72176687)(303.45820313,127.41082937)
\curveto(302.70351563,128.10809552)(302.32617188,129.07196192)(302.32617188,130.29832937)
\curveto(302.32617188,131.53289942)(302.70351563,132.49676687)(303.45820313,133.18582937)
\curveto(304.22109427,133.87489187)(305.27929688,134.21942312)(306.62460938,134.21942312)
\curveto(307.01835937,134.21942312)(307.39980443,134.17840802)(307.77304688,134.10457937)
\curveto(308.15449193,134.03895437)(308.53593698,133.94051687)(308.92148438,133.80926687)
\closepath
\moveto(308.92148438,133.80926687)
}
}
{
\newrgbcolor{curcolor}{0 0 0}
\pscustom[linestyle=none,fillstyle=solid,fillcolor=curcolor]
{
\newpath
\moveto(317.02207031,133.80926687)
\lineto(317.02207031,131.85692312)
\curveto(316.69394531,132.07430567)(316.36582031,132.23836817)(316.03769531,132.34911062)
\curveto(315.70957031,132.45575072)(315.36914036,132.51317312)(315.02050781,132.51317312)
\curveto(314.35195286,132.51317312)(313.83105521,132.31629812)(313.46191406,131.92254812)
\curveto(313.08867161,131.53700072)(312.90410156,130.99559552)(312.90410156,130.29832937)
\curveto(312.90410156,129.59696192)(313.08867161,129.05145437)(313.46191406,128.65770437)
\curveto(313.83105521,128.27215802)(314.35195286,128.08348562)(315.02050781,128.08348562)
\curveto(315.38964791,128.08348562)(315.74648411,128.13680567)(316.08691406,128.24754812)
\curveto(316.42324271,128.36649317)(316.73496041,128.53055567)(317.02207031,128.73973562)
\lineto(317.02207031,126.78739187)
\curveto(316.63652291,126.65614187)(316.25507786,126.55770437)(315.87363281,126.49207937)
\curveto(315.48808541,126.41825072)(315.10664036,126.37723562)(314.72519531,126.37723562)
\curveto(313.37988281,126.37723562)(312.32168021,126.72176687)(311.55878906,127.41082937)
\curveto(310.80410156,128.10809552)(310.42675781,129.07196192)(310.42675781,130.29832937)
\curveto(310.42675781,131.53289942)(310.80410156,132.49676687)(311.55878906,133.18582937)
\curveto(312.32168021,133.87489187)(313.37988281,134.21942312)(314.72519531,134.21942312)
\curveto(315.11894531,134.21942312)(315.50039036,134.17840802)(315.87363281,134.10457937)
\curveto(316.25507786,134.03895437)(316.63652291,133.94051687)(317.02207031,133.80926687)
\closepath
\moveto(317.02207031,133.80926687)
}
}
{
\newrgbcolor{curcolor}{0 0 0}
\pscustom[linestyle=none,fillstyle=solid,fillcolor=curcolor]
{
\newpath
\moveto(326.53359375,130.33114187)
\lineto(326.53359375,129.64207937)
\lineto(320.95546875,129.64207937)
\curveto(321.0087888,129.08426687)(321.20976615,128.66180567)(321.5625,128.37879812)
\curveto(321.91113255,128.10399317)(322.39921875,127.96864187)(323.02265625,127.96864187)
\curveto(323.51484375,127.96864187)(324.0234375,128.03836817)(324.5484375,128.18192312)
\curveto(325.08164115,128.33368067)(325.63125,128.56336817)(326.1890625,128.87098562)
\lineto(326.1890625,127.03348562)
\curveto(325.61894505,126.81200072)(325.04882865,126.64793875)(324.4828125,126.54129812)
\curveto(323.925,126.43465802)(323.35898385,126.37723562)(322.79296875,126.37723562)
\curveto(321.44765625,126.37723562)(320.39765625,126.72176687)(319.64296875,127.41082937)
\curveto(318.89648385,128.09989187)(318.52734375,129.05965802)(318.52734375,130.29832937)
\curveto(318.52734375,131.51239187)(318.89238255,132.46805567)(319.6265625,133.16942312)
\curveto(320.35664115,133.86668822)(321.3697263,134.21942312)(322.66171875,134.21942312)
\curveto(323.8306638,134.21942312)(324.76582005,133.86258692)(325.4671875,133.15301687)
\curveto(326.17675755,132.45164942)(326.53359375,131.51239187)(326.53359375,130.33114187)
\closepath
\moveto(324.0890625,131.11864187)
\curveto(324.0890625,131.57801687)(323.94960885,131.94305567)(323.67890625,132.21786062)
\curveto(323.41640625,132.48856322)(323.071875,132.62801687)(322.6453125,132.62801687)
\curveto(322.1859375,132.62801687)(321.80859375,132.49676687)(321.51328125,132.23426687)
\curveto(321.22617135,131.97997052)(321.05390625,131.61082937)(320.98828125,131.11864187)
\closepath
\moveto(324.0890625,131.11864187)
}
}
{
\newrgbcolor{curcolor}{0 0 0}
\pscustom[linestyle=none,fillstyle=solid,fillcolor=curcolor]
{
\newpath
\moveto(330.71308594,127.65692312)
\lineto(330.71308594,123.73582937)
\lineto(328.33417969,123.73582937)
\lineto(328.33417969,134.03895437)
\lineto(330.71308594,134.03895437)
\lineto(330.71308594,132.93973562)
\curveto(331.04121094,133.37450072)(331.40624974,133.69852442)(331.81230469,133.90770437)
\curveto(332.21425729,134.11278302)(332.67363229,134.21942312)(333.19042969,134.21942312)
\curveto(334.10917969,134.21942312)(334.86386719,133.85028302)(335.45449219,133.12020437)
\curveto(336.04511719,132.38602442)(336.34042969,131.44676687)(336.34042969,130.29832937)
\curveto(336.34042969,129.14989187)(336.04511719,128.20653302)(335.45449219,127.47645437)
\curveto(334.86386719,126.74227442)(334.10917969,126.37723562)(333.19042969,126.37723562)
\curveto(332.67363229,126.37723562)(332.21425729,126.48387625)(331.81230469,126.68895437)
\curveto(331.40624974,126.89403302)(331.04121094,127.21805567)(330.71308594,127.65692312)
\closepath
\moveto(332.30449219,132.48036062)
\curveto(331.78769479,132.48036062)(331.39394479,132.29168822)(331.12324219,131.92254812)
\curveto(330.84843724,131.54930567)(330.71308594,131.00789942)(330.71308594,130.29832937)
\curveto(330.71308594,129.59696192)(330.84843724,129.05555567)(331.12324219,128.67411062)
\curveto(331.39394479,128.30086817)(331.78769479,128.11629812)(332.30449219,128.11629812)
\curveto(332.81718724,128.11629812)(333.20683594,128.30086817)(333.46933594,128.67411062)
\curveto(333.74003959,129.04325072)(333.87949219,129.58465802)(333.87949219,130.29832937)
\curveto(333.87949219,131.00789942)(333.74003959,131.54930567)(333.46933594,131.92254812)
\curveto(333.20683594,132.29168822)(332.81718724,132.48036062)(332.30449219,132.48036062)
\closepath
\moveto(332.30449219,132.48036062)
}
}
{
\newrgbcolor{curcolor}{0 0 0}
\pscustom[linestyle=none,fillstyle=solid,fillcolor=curcolor]
{
\newpath
\moveto(340.70449219,136.15536062)
\lineto(340.70449219,134.03895437)
\lineto(343.16542969,134.03895437)
\lineto(343.16542969,132.33270437)
\lineto(340.70449219,132.33270437)
\lineto(340.70449219,129.16629812)
\curveto(340.70449219,128.81356322)(340.77011719,128.57977442)(340.90136719,128.46082937)
\curveto(341.04082084,128.33778302)(341.31562474,128.28036062)(341.72167969,128.28036062)
\lineto(342.95214844,128.28036062)
\lineto(342.95214844,126.57411062)
\lineto(340.90136719,126.57411062)
\curveto(339.95800834,126.57411062)(339.29355469,126.77098562)(338.89980469,127.16473562)
\curveto(338.50605469,127.55848562)(338.30917969,128.22293822)(338.30917969,129.16629812)
\lineto(338.30917969,132.33270437)
\lineto(337.12792969,132.33270437)
\lineto(337.12792969,134.03895437)
\lineto(338.30917969,134.03895437)
\lineto(338.30917969,136.15536062)
\closepath
\moveto(340.70449219,136.15536062)
}
}
{
\newrgbcolor{curcolor}{0 0 0}
\pscustom[linestyle=none,fillstyle=solid,fillcolor=curcolor]
{
\newpath
\moveto(347.96425781,129.93739187)
\curveto(347.45976536,129.93739187)(347.08242161,129.84715802)(346.83222656,129.67489187)
\curveto(346.57793021,129.51082937)(346.45488281,129.26473562)(346.45488281,128.93661062)
\curveto(346.45488281,128.62899317)(346.55332031,128.38700072)(346.75019531,128.21473562)
\curveto(346.95527291,128.05067312)(347.24238281,127.96864187)(347.60332031,127.96864187)
\curveto(348.05039036,127.96864187)(348.42773411,128.12450072)(348.73535156,128.44442312)
\curveto(349.03886771,128.77254812)(349.19472656,129.17450072)(349.19472656,129.65848562)
\lineto(349.19472656,129.93739187)
\closepath
\moveto(351.60644531,130.83973562)
\lineto(351.60644531,126.57411062)
\lineto(349.19472656,126.57411062)
\lineto(349.19472656,127.67332937)
\curveto(348.87480521,127.22215802)(348.51386771,126.89403302)(348.11191406,126.68895437)
\curveto(347.71816406,126.48387625)(347.23418021,126.37723562)(346.66816406,126.37723562)
\curveto(345.88886771,126.37723562)(345.26132786,126.59872)(344.78144531,127.04989187)
\curveto(344.29746041,127.50926687)(344.05957031,128.09989187)(344.05957031,128.82176687)
\curveto(344.05957031,129.69539942)(344.35898411,130.33524317)(344.96191406,130.74129812)
\curveto(345.56074271,131.14325072)(346.50820286,131.34832937)(347.80019531,131.34832937)
\lineto(349.19472656,131.34832937)
\lineto(349.19472656,131.54520437)
\curveto(349.19472656,131.91434552)(349.04707031,132.18914942)(348.75175781,132.36551687)
\curveto(348.45644531,132.53778302)(347.98886771,132.62801687)(347.35722656,132.62801687)
\curveto(346.85273411,132.62801687)(346.37695286,132.57879812)(345.92988281,132.48036062)
\curveto(345.49101536,132.38192312)(345.08496041,132.22606322)(344.71582031,132.02098562)
\lineto(344.71582031,133.84207937)
\curveto(345.21621041,133.96102442)(345.72480521,134.05536062)(346.24160156,134.12098562)
\curveto(346.76660156,134.18661062)(347.28339791,134.21942312)(347.80019531,134.21942312)
\curveto(349.14550781,134.21942312)(350.11347656,133.94872052)(350.70410156,133.41551687)
\curveto(351.30293021,132.89051687)(351.60644531,132.02918822)(351.60644531,130.83973562)
\closepath
\moveto(351.60644531,130.83973562)
}
}
{
\newrgbcolor{curcolor}{0 0 0}
\pscustom[linestyle=none,fillstyle=solid,fillcolor=curcolor]
{
\newpath
\moveto(361.32304687,131.11864187)
\lineto(361.32304687,126.57411062)
\lineto(358.92773438,126.57411062)
\lineto(358.92773438,130.05223562)
\curveto(358.92773438,130.69618067)(358.91132813,131.13914942)(358.87851563,131.38114187)
\curveto(358.85390573,131.63133692)(358.80468803,131.81180567)(358.73085938,131.92254812)
\curveto(358.64062553,132.07430567)(358.51757812,132.19325072)(358.35351562,132.28348562)
\curveto(358.19765573,132.36961817)(358.01718802,132.41473562)(357.81210938,132.41473562)
\curveto(357.30761692,132.41473562)(356.91386693,132.22196192)(356.63085938,131.84051687)
\curveto(356.34375053,131.45497052)(356.20429688,130.92586817)(356.20429688,130.24911062)
\lineto(356.20429688,126.57411062)
\lineto(353.82539063,126.57411062)
\lineto(353.82539063,134.03895437)
\lineto(356.20429688,134.03895437)
\lineto(356.20429688,132.93973562)
\curveto(356.56523437,133.37450072)(356.94667943,133.69852442)(357.35273438,133.90770437)
\curveto(357.75468803,134.11278302)(358.20585937,134.21942312)(358.69804687,134.21942312)
\curveto(359.55937553,134.21942312)(360.21152318,133.94872052)(360.65039062,133.41551687)
\curveto(361.09746068,132.89051687)(361.32304687,132.12352442)(361.32304687,131.11864187)
\closepath
\moveto(361.32304687,131.11864187)
}
}
{
\newrgbcolor{curcolor}{0 0 0}
\pscustom[linestyle=none,fillstyle=solid,fillcolor=curcolor]
{
\newpath
\moveto(369.58359375,133.80926687)
\lineto(369.58359375,131.85692312)
\curveto(369.25546875,132.07430567)(368.92734375,132.23836817)(368.59921875,132.34911062)
\curveto(368.27109375,132.45575072)(367.9306638,132.51317312)(367.58203125,132.51317312)
\curveto(366.9134763,132.51317312)(366.39257865,132.31629812)(366.0234375,131.92254812)
\curveto(365.65019505,131.53700072)(365.465625,130.99559552)(365.465625,130.29832937)
\curveto(365.465625,129.59696192)(365.65019505,129.05145437)(366.0234375,128.65770437)
\curveto(366.39257865,128.27215802)(366.9134763,128.08348562)(367.58203125,128.08348562)
\curveto(367.95117135,128.08348562)(368.30800755,128.13680567)(368.6484375,128.24754812)
\curveto(368.98476615,128.36649317)(369.29648385,128.53055567)(369.58359375,128.73973562)
\lineto(369.58359375,126.78739187)
\curveto(369.19804635,126.65614187)(368.8166013,126.55770437)(368.43515625,126.49207937)
\curveto(368.04960885,126.41825072)(367.6681638,126.37723562)(367.28671875,126.37723562)
\curveto(365.94140625,126.37723562)(364.88320365,126.72176687)(364.1203125,127.41082937)
\curveto(363.365625,128.10809552)(362.98828125,129.07196192)(362.98828125,130.29832937)
\curveto(362.98828125,131.53289942)(363.365625,132.49676687)(364.1203125,133.18582937)
\curveto(364.88320365,133.87489187)(365.94140625,134.21942312)(367.28671875,134.21942312)
\curveto(367.68046875,134.21942312)(368.0619138,134.17840802)(368.43515625,134.10457937)
\curveto(368.8166013,134.03895437)(369.19804635,133.94051687)(369.58359375,133.80926687)
\closepath
\moveto(369.58359375,133.80926687)
}
}
{
\newrgbcolor{curcolor}{0 0 0}
\pscustom[linestyle=none,fillstyle=solid,fillcolor=curcolor]
{
\newpath
\moveto(379.09511719,130.33114187)
\lineto(379.09511719,129.64207937)
\lineto(373.51699219,129.64207937)
\curveto(373.57031224,129.08426687)(373.77128959,128.66180567)(374.12402344,128.37879812)
\curveto(374.47265599,128.10399317)(374.96074219,127.96864187)(375.58417969,127.96864187)
\curveto(376.07636719,127.96864187)(376.58496094,128.03836817)(377.10996094,128.18192312)
\curveto(377.64316459,128.33368067)(378.19277344,128.56336817)(378.75058594,128.87098562)
\lineto(378.75058594,127.03348562)
\curveto(378.18046849,126.81200072)(377.61035209,126.64793875)(377.04433594,126.54129812)
\curveto(376.48652344,126.43465802)(375.92050729,126.37723562)(375.35449219,126.37723562)
\curveto(374.00917969,126.37723562)(372.95917969,126.72176687)(372.20449219,127.41082937)
\curveto(371.45800729,128.09989187)(371.08886719,129.05965802)(371.08886719,130.29832937)
\curveto(371.08886719,131.51239187)(371.45390599,132.46805567)(372.18808594,133.16942312)
\curveto(372.91816459,133.86668822)(373.93124974,134.21942312)(375.22324219,134.21942312)
\curveto(376.39218724,134.21942312)(377.32734349,133.86258692)(378.02871094,133.15301687)
\curveto(378.73828099,132.45164942)(379.09511719,131.51239187)(379.09511719,130.33114187)
\closepath
\moveto(376.65058594,131.11864187)
\curveto(376.65058594,131.57801687)(376.51113229,131.94305567)(376.24042969,132.21786062)
\curveto(375.97792969,132.48856322)(375.63339844,132.62801687)(375.20683594,132.62801687)
\curveto(374.74746094,132.62801687)(374.37011719,132.49676687)(374.07480469,132.23426687)
\curveto(373.78769479,131.97997052)(373.61542969,131.61082937)(373.54980469,131.11864187)
\closepath
\moveto(376.65058594,131.11864187)
}
}
{
\newrgbcolor{curcolor}{0 0 0}
\pscustom[linewidth=2.10000003,linecolor=curcolor]
{
\newpath
\moveto(519.35625,265.20692575)
\curveto(530.647845,212.44442575)(445.499418,131.55750775)(412.051164,130.87664575)
}
}
{
\newrgbcolor{curcolor}{0 0 0}
\pscustom[linestyle=none,fillstyle=solid,fillcolor=curcolor]
{
\newpath
\moveto(404.176164,130.71668875)
\lineto(414.565431,136.17996475)
\lineto(412.051164,130.87664575)
\lineto(414.778707,125.67996475)
\closepath
\moveto(404.176164,130.71668875)
}
}
{
\newrgbcolor{curcolor}{0 0 0}
\pscustom[linewidth=2.10000003,linecolor=curcolor]
{
\newpath
\moveto(404.176164,130.71668875)
\lineto(414.565431,136.17996475)
\lineto(412.051164,130.87664575)
\lineto(414.778707,125.67996475)
\closepath
\moveto(404.176164,130.71668875)
}
}
{
\newrgbcolor{curcolor}{0 0 0}
\pscustom[linewidth=2.10000003,linecolor=curcolor]
{
\newpath
\moveto(330.007608,40.34285275)
\lineto(480.526767,40.34285275)
\lineto(480.526767,0.44285275)
\lineto(330.007608,0.44285275)
\closepath
\moveto(330.007608,40.34285275)
}
}
{
\newrgbcolor{curcolor}{0 0 0}
\pscustom[linestyle=none,fillstyle=solid,fillcolor=curcolor]
{
\newpath
\moveto(365.465625,26.258095)
\lineto(374.63671875,26.258095)
\lineto(374.63671875,24.30575125)
\lineto(371.3390625,24.30575125)
\lineto(371.3390625,16.29950125)
\lineto(368.7796875,16.29950125)
\lineto(368.7796875,24.30575125)
\lineto(365.465625,24.30575125)
\closepath
\moveto(365.465625,26.258095)
}
}
{
\newrgbcolor{curcolor}{0 0 0}
\pscustom[linestyle=none,fillstyle=solid,fillcolor=curcolor]
{
\newpath
\moveto(383.35664063,20.8440325)
\lineto(383.35664063,16.29950125)
\lineto(380.96132813,16.29950125)
\lineto(380.96132813,19.76122)
\curveto(380.96132813,20.41747)(380.94492188,20.86454005)(380.91210938,21.1065325)
\curveto(380.88749948,21.35672755)(380.83828178,21.5371963)(380.76445313,21.64793875)
\curveto(380.67421928,21.7996963)(380.55117188,21.91864135)(380.38710938,22.00887625)
\curveto(380.23124948,22.0950088)(380.05078178,22.14012625)(379.84570313,22.14012625)
\curveto(379.34121068,22.14012625)(378.94746068,21.94735255)(378.66445313,21.5659075)
\curveto(378.37734428,21.18036115)(378.23789063,20.6512588)(378.23789063,19.97450125)
\lineto(378.23789063,16.29950125)
\lineto(375.85898438,16.29950125)
\lineto(375.85898438,26.66825125)
\lineto(378.23789063,26.66825125)
\lineto(378.23789063,22.66512625)
\curveto(378.59882813,23.09989135)(378.98027318,23.42391505)(379.38632813,23.633095)
\curveto(379.78828178,23.83817365)(380.23945313,23.94481375)(380.73164063,23.94481375)
\curveto(381.59296928,23.94481375)(382.24511693,23.67411115)(382.68398438,23.1409075)
\curveto(383.13105443,22.6159075)(383.35664063,21.84891505)(383.35664063,20.8440325)
\closepath
\moveto(383.35664063,20.8440325)
}
}
{
\newrgbcolor{curcolor}{0 0 0}
\pscustom[linestyle=none,fillstyle=solid,fillcolor=curcolor]
{
\newpath
\moveto(388.9265625,19.6627825)
\curveto(388.42207005,19.6627825)(388.0447263,19.57254865)(387.79453125,19.4002825)
\curveto(387.5402349,19.23622)(387.4171875,18.99012625)(387.4171875,18.66200125)
\curveto(387.4171875,18.3543838)(387.515625,18.11239135)(387.7125,17.94012625)
\curveto(387.9175776,17.77606375)(388.2046875,17.6940325)(388.565625,17.6940325)
\curveto(389.01269505,17.6940325)(389.3900388,17.84989135)(389.69765625,18.16981375)
\curveto(390.0011724,18.49793875)(390.15703125,18.89989135)(390.15703125,19.38387625)
\lineto(390.15703125,19.6627825)
\closepath
\moveto(392.56875,20.56512625)
\lineto(392.56875,16.29950125)
\lineto(390.15703125,16.29950125)
\lineto(390.15703125,17.39872)
\curveto(389.8371099,16.94754865)(389.4761724,16.61942365)(389.07421875,16.414345)
\curveto(388.68046875,16.20926687)(388.1964849,16.10262625)(387.63046875,16.10262625)
\curveto(386.8511724,16.10262625)(386.22363255,16.32411062)(385.74375,16.7752825)
\curveto(385.2597651,17.2346575)(385.021875,17.8252825)(385.021875,18.5471575)
\curveto(385.021875,19.42079005)(385.3212888,20.0606338)(385.92421875,20.46668875)
\curveto(386.5230474,20.86864135)(387.47050755,21.07372)(388.7625,21.07372)
\lineto(390.15703125,21.07372)
\lineto(390.15703125,21.270595)
\curveto(390.15703125,21.63973615)(390.009375,21.91454005)(389.7140625,22.0909075)
\curveto(389.41875,22.26317365)(388.9511724,22.3534075)(388.31953125,22.3534075)
\curveto(387.8150388,22.3534075)(387.33925755,22.30418875)(386.8921875,22.20575125)
\curveto(386.45332005,22.10731375)(386.0472651,21.95145385)(385.678125,21.74637625)
\lineto(385.678125,23.56747)
\curveto(386.1785151,23.68641505)(386.6871099,23.78075125)(387.20390625,23.84637625)
\curveto(387.72890625,23.91200125)(388.2457026,23.94481375)(388.7625,23.94481375)
\curveto(390.1078125,23.94481375)(391.07578125,23.67411115)(391.66640625,23.1409075)
\curveto(392.2652349,22.6159075)(392.56875,21.75457885)(392.56875,20.56512625)
\closepath
\moveto(392.56875,20.56512625)
}
}
{
\newrgbcolor{curcolor}{0 0 0}
\pscustom[linestyle=none,fillstyle=solid,fillcolor=curcolor]
{
\newpath
\moveto(402.28535156,20.8440325)
\lineto(402.28535156,16.29950125)
\lineto(399.89003906,16.29950125)
\lineto(399.89003906,19.77762625)
\curveto(399.89003906,20.4215713)(399.87363281,20.86454005)(399.84082031,21.1065325)
\curveto(399.81621041,21.35672755)(399.76699271,21.5371963)(399.69316406,21.64793875)
\curveto(399.60293021,21.7996963)(399.47988281,21.91864135)(399.31582031,22.00887625)
\curveto(399.15996041,22.0950088)(398.97949271,22.14012625)(398.77441406,22.14012625)
\curveto(398.26992161,22.14012625)(397.87617161,21.94735255)(397.59316406,21.5659075)
\curveto(397.30605521,21.18036115)(397.16660156,20.6512588)(397.16660156,19.97450125)
\lineto(397.16660156,16.29950125)
\lineto(394.78769531,16.29950125)
\lineto(394.78769531,23.764345)
\lineto(397.16660156,23.764345)
\lineto(397.16660156,22.66512625)
\curveto(397.52753906,23.09989135)(397.90898411,23.42391505)(398.31503906,23.633095)
\curveto(398.71699271,23.83817365)(399.16816406,23.94481375)(399.66035156,23.94481375)
\curveto(400.52168021,23.94481375)(401.17382786,23.67411115)(401.61269531,23.1409075)
\curveto(402.05976536,22.6159075)(402.28535156,21.84891505)(402.28535156,20.8440325)
\closepath
\moveto(402.28535156,20.8440325)
}
}
{
\newrgbcolor{curcolor}{0 0 0}
\pscustom[linestyle=none,fillstyle=solid,fillcolor=curcolor]
{
\newpath
\moveto(404.50839844,26.66825125)
\lineto(406.88730469,26.66825125)
\lineto(406.88730469,21.02450125)
\lineto(409.64355469,23.764345)
\lineto(412.41621094,23.764345)
\lineto(408.75761719,20.33543875)
\lineto(412.69511719,16.29950125)
\lineto(409.80761719,16.29950125)
\lineto(406.88730469,19.41668875)
\lineto(406.88730469,16.29950125)
\lineto(404.50839844,16.29950125)
\closepath
\moveto(404.50839844,26.66825125)
}
}
{
\newrgbcolor{curcolor}{0 0 0}
\pscustom[linestyle=none,fillstyle=solid,fillcolor=curcolor]
{
\newpath
\moveto(417.07148438,26.258095)
\lineto(419.87695313,26.258095)
\lineto(422.14101563,22.69793875)
\lineto(424.40507813,26.258095)
\lineto(427.22695313,26.258095)
\lineto(423.42070313,20.49950125)
\lineto(423.42070313,16.29950125)
\lineto(420.86132813,16.29950125)
\lineto(420.86132813,20.49950125)
\closepath
\moveto(417.07148438,26.258095)
}
}
{
\newrgbcolor{curcolor}{0 0 0}
\pscustom[linestyle=none,fillstyle=solid,fillcolor=curcolor]
{
\newpath
\moveto(430.73378906,22.23856375)
\curveto(430.20878906,22.23856375)(429.80273411,22.04579005)(429.51972656,21.664345)
\curveto(429.24492161,21.29110255)(429.10957031,20.745595)(429.10957031,20.02372)
\curveto(429.10957031,19.31004865)(429.24492161,18.76454005)(429.51972656,18.383095)
\curveto(429.80273411,17.99754865)(430.20878906,17.80887625)(430.73378906,17.80887625)
\curveto(431.25878906,17.80887625)(431.65664036,17.99754865)(431.93144531,18.383095)
\curveto(432.20214896,18.76454005)(432.34160156,19.31004865)(432.34160156,20.02372)
\curveto(432.34160156,20.745595)(432.20214896,21.29110255)(431.93144531,21.664345)
\curveto(431.65664036,22.04579005)(431.25878906,22.23856375)(430.73378906,22.23856375)
\closepath
\moveto(430.73378906,23.94481375)
\curveto(432.02167916,23.94481375)(433.03066406,23.59207885)(433.75253906,22.89481375)
\curveto(434.47441406,22.20575125)(434.83535156,21.24598615)(434.83535156,20.02372)
\curveto(434.83535156,18.79735255)(434.47441406,17.83348615)(433.75253906,17.13622)
\curveto(433.03066406,16.4471575)(432.02167916,16.10262625)(430.73378906,16.10262625)
\curveto(429.45410156,16.10262625)(428.44511666,16.4471575)(427.71503906,17.13622)
\curveto(426.99316406,17.83348615)(426.63222656,18.79735255)(426.63222656,20.02372)
\curveto(426.63222656,21.24598615)(426.99316406,22.20575125)(427.71503906,22.89481375)
\curveto(428.44511666,23.59207885)(429.45410156,23.94481375)(430.73378906,23.94481375)
\closepath
\moveto(430.73378906,23.94481375)
}
}
{
\newrgbcolor{curcolor}{0 0 0}
\pscustom[linestyle=none,fillstyle=solid,fillcolor=curcolor]
{
\newpath
\moveto(436.48007813,19.2034075)
\lineto(436.48007813,23.764345)
\lineto(438.87539063,23.764345)
\lineto(438.87539063,23.02606375)
\curveto(438.87539063,22.6200088)(438.86718698,22.11141505)(438.85898438,21.5002825)
\lineto(438.85898438,20.26981375)
\curveto(438.85898438,19.6668838)(438.87539063,19.23622)(438.90820313,18.97372)
\curveto(438.94101563,18.71122)(438.99433568,18.5184463)(439.07226563,18.39950125)
\curveto(439.17070313,18.24364135)(439.29374948,18.1246963)(439.44960938,18.03856375)
\curveto(439.61367188,17.94832885)(439.79824193,17.90731375)(440.00742188,17.90731375)
\curveto(440.49960938,17.90731375)(440.88515678,18.09598615)(441.17226563,18.4815325)
\curveto(441.45527318,18.86297755)(441.59882813,19.39207885)(441.59882813,20.07293875)
\lineto(441.59882813,23.764345)
\lineto(443.97773438,23.764345)
\lineto(443.97773438,16.29950125)
\lineto(441.59882813,16.29950125)
\lineto(441.59882813,17.38231375)
\curveto(441.23789063,16.9434463)(440.85234428,16.61942365)(440.45039063,16.414345)
\curveto(440.05664063,16.20926687)(439.61367188,16.10262625)(439.12148438,16.10262625)
\curveto(438.25605443,16.10262625)(437.59980443,16.36512625)(437.15273438,16.89012625)
\curveto(436.70156198,17.42332885)(436.48007813,18.19442365)(436.48007813,19.2034075)
\closepath
\moveto(436.48007813,19.2034075)
}
}
{
\newrgbcolor{curcolor}{0 0 0}
\pscustom[linewidth=2.10000003,linecolor=curcolor]
{
\newpath
\moveto(75.956832,102.80145175)
\curveto(89.586336,15.75399475)(277.761918,9.79442575)(319.9875,18.36259375)
}
}
{
\newrgbcolor{curcolor}{0 0 0}
\pscustom[linestyle=none,fillstyle=solid,fillcolor=curcolor]
{
\newpath
\moveto(327.706638,19.92938275)
\lineto(318.457608,12.69423175)
\lineto(319.987458,18.36259375)
\lineto(316.369872,22.98505075)
\closepath
\moveto(327.706638,19.92938275)
}
}
{
\newrgbcolor{curcolor}{0 0 0}
\pscustom[linewidth=2.10000003,linecolor=curcolor]
{
\newpath
\moveto(327.706638,19.92938275)
\lineto(318.457608,12.69423175)
\lineto(319.987458,18.36259375)
\lineto(316.369872,22.98505075)
\closepath
\moveto(327.706638,19.92938275)
}
}
{
\newrgbcolor{curcolor}{0 0 0}
\pscustom[linewidth=2.10000003,linecolor=curcolor]
{
\newpath
\moveto(269.070711,130.66746475)
\curveto(192.703707,127.09910275)(113.978319,202.44071875)(113.945517,255.53543875)
}
}
{
\newrgbcolor{curcolor}{0 0 0}
\pscustom[linestyle=none,fillstyle=solid,fillcolor=curcolor]
{
\newpath
\moveto(113.941401,263.41043875)
\lineto(119.195517,252.91043875)
\lineto(113.945517,255.53543875)
\lineto(108.695517,252.90634375)
\closepath
\moveto(113.941401,263.41043875)
}
}
{
\newrgbcolor{curcolor}{0 0 0}
\pscustom[linewidth=2.10000003,linecolor=curcolor]
{
\newpath
\moveto(113.941401,263.41043875)
\lineto(119.195517,252.91043875)
\lineto(113.945517,255.53543875)
\lineto(108.695517,252.90634375)
\closepath
\moveto(113.941401,263.41043875)
}
}
{
\newrgbcolor{curcolor}{0 0 0}
\pscustom[linewidth=2.10000003,linecolor=curcolor]
{
\newpath
\moveto(222.571293,304.67216275)
\lineto(373.090431,304.67216275)
\lineto(373.090431,264.77216275)
\lineto(222.571293,264.77216275)
\closepath
\moveto(222.571293,304.67216275)
}
}
{
\newrgbcolor{curcolor}{0 0 0}
\pscustom[linestyle=none,fillstyle=solid,fillcolor=curcolor]
{
\newpath
\moveto(259.96933594,290.58739187)
\lineto(262.84042969,290.58739187)
\lineto(266.46621094,283.76239187)
\lineto(266.46621094,290.58739187)
\lineto(268.89433594,290.58739187)
\lineto(268.89433594,280.62879812)
\lineto(266.02324219,280.62879812)
\lineto(262.41386719,287.45379812)
\lineto(262.41386719,280.62879812)
\lineto(259.96933594,280.62879812)
\closepath
\moveto(259.96933594,290.58739187)
}
}
{
\newrgbcolor{curcolor}{0 0 0}
\pscustom[linestyle=none,fillstyle=solid,fillcolor=curcolor]
{
\newpath
\moveto(278.7421875,284.38582937)
\lineto(278.7421875,283.69676687)
\lineto(273.1640625,283.69676687)
\curveto(273.21738255,283.13895437)(273.4183599,282.71649317)(273.77109375,282.43348562)
\curveto(274.1197263,282.15868067)(274.6078125,282.02332937)(275.23125,282.02332937)
\curveto(275.7234375,282.02332937)(276.23203125,282.09305567)(276.75703125,282.23661062)
\curveto(277.2902349,282.38836817)(277.83984375,282.61805567)(278.39765625,282.92567312)
\lineto(278.39765625,281.08817312)
\curveto(277.8275388,280.86668822)(277.2574224,280.70262625)(276.69140625,280.59598562)
\curveto(276.13359375,280.48934552)(275.5675776,280.43192312)(275.0015625,280.43192312)
\curveto(273.65625,280.43192312)(272.60625,280.77645437)(271.8515625,281.46551687)
\curveto(271.1050776,282.15457937)(270.7359375,283.11434552)(270.7359375,284.35301687)
\curveto(270.7359375,285.56707937)(271.1009763,286.52274317)(271.83515625,287.22411062)
\curveto(272.5652349,287.92137572)(273.57832005,288.27411062)(274.8703125,288.27411062)
\curveto(276.03925755,288.27411062)(276.9744138,287.91727442)(277.67578125,287.20770437)
\curveto(278.3853513,286.50633692)(278.7421875,285.56707937)(278.7421875,284.38582937)
\closepath
\moveto(276.29765625,285.17332937)
\curveto(276.29765625,285.63270437)(276.1582026,285.99774317)(275.8875,286.27254812)
\curveto(275.625,286.54325072)(275.28046875,286.68270437)(274.85390625,286.68270437)
\curveto(274.39453125,286.68270437)(274.0171875,286.55145437)(273.721875,286.28895437)
\curveto(273.4347651,286.03465802)(273.2625,285.66551687)(273.196875,285.17332937)
\closepath
\moveto(276.29765625,285.17332937)
}
}
{
\newrgbcolor{curcolor}{0 0 0}
\pscustom[linestyle=none,fillstyle=solid,fillcolor=curcolor]
{
\newpath
\moveto(282.42949219,284.45145437)
\lineto(279.73886719,288.09364187)
\lineto(282.26542969,288.09364187)
\lineto(283.79121094,285.87879812)
\lineto(285.33339844,288.09364187)
\lineto(287.85996094,288.09364187)
\lineto(285.16933594,284.46786062)
\lineto(287.99121094,280.62879812)
\lineto(285.46464844,280.62879812)
\lineto(283.79121094,282.99129812)
\lineto(282.13417969,280.62879812)
\lineto(279.60761719,280.62879812)
\closepath
\moveto(282.42949219,284.45145437)
}
}
{
\newrgbcolor{curcolor}{0 0 0}
\pscustom[linestyle=none,fillstyle=solid,fillcolor=curcolor]
{
\newpath
\moveto(291.94921875,290.21004812)
\lineto(291.94921875,288.09364187)
\lineto(294.41015625,288.09364187)
\lineto(294.41015625,286.38739187)
\lineto(291.94921875,286.38739187)
\lineto(291.94921875,283.22098562)
\curveto(291.94921875,282.86825072)(292.01484375,282.63446192)(292.14609375,282.51551687)
\curveto(292.2855474,282.39247052)(292.5603513,282.33504812)(292.96640625,282.33504812)
\lineto(294.196875,282.33504812)
\lineto(294.196875,280.62879812)
\lineto(292.14609375,280.62879812)
\curveto(291.2027349,280.62879812)(290.53828125,280.82567312)(290.14453125,281.21942312)
\curveto(289.75078125,281.61317312)(289.55390625,282.27762572)(289.55390625,283.22098562)
\lineto(289.55390625,286.38739187)
\lineto(288.37265625,286.38739187)
\lineto(288.37265625,288.09364187)
\lineto(289.55390625,288.09364187)
\lineto(289.55390625,290.21004812)
\closepath
\moveto(291.94921875,290.21004812)
}
}
{
\newrgbcolor{curcolor}{0 0 0}
\pscustom[linestyle=none,fillstyle=solid,fillcolor=curcolor]
{
\newpath
\moveto(299.53710938,290.58739187)
\lineto(308.70820313,290.58739187)
\lineto(308.70820313,288.63504812)
\lineto(305.41054687,288.63504812)
\lineto(305.41054687,280.62879812)
\lineto(302.85117187,280.62879812)
\lineto(302.85117187,288.63504812)
\lineto(299.53710938,288.63504812)
\closepath
\moveto(299.53710938,290.58739187)
}
}
{
\newrgbcolor{curcolor}{0 0 0}
\pscustom[linestyle=none,fillstyle=solid,fillcolor=curcolor]
{
\newpath
\moveto(309.93046875,288.09364187)
\lineto(312.309375,288.09364187)
\lineto(312.309375,280.62879812)
\lineto(309.93046875,280.62879812)
\closepath
\moveto(309.93046875,290.99754812)
\lineto(312.309375,290.99754812)
\lineto(312.309375,289.06161062)
\lineto(309.93046875,289.06161062)
\closepath
\moveto(309.93046875,290.99754812)
}
}
{
\newrgbcolor{curcolor}{0 0 0}
\pscustom[linestyle=none,fillstyle=solid,fillcolor=curcolor]
{
\newpath
\moveto(321.5296875,286.84676687)
\curveto(321.825,287.31434552)(322.1777349,287.67118067)(322.59609375,287.91317312)
\curveto(323.0103513,288.15106322)(323.4697263,288.27411062)(323.97421875,288.27411062)
\curveto(324.82734375,288.27411062)(325.4753901,288.00340802)(325.9265625,287.47020437)
\curveto(326.37363255,286.94520437)(326.59921875,286.17821192)(326.59921875,285.17332937)
\lineto(326.59921875,280.62879812)
\lineto(324.20390625,280.62879812)
\lineto(324.20390625,284.51707937)
\lineto(324.20390625,284.69754812)
\curveto(324.2121099,284.76317312)(324.2203125,284.84930567)(324.2203125,284.96004812)
\curveto(324.2203125,285.49325072)(324.13828125,285.87879812)(323.97421875,286.10848562)
\curveto(323.8183599,286.34637572)(323.5722651,286.46942312)(323.2359375,286.46942312)
\curveto(322.7765625,286.46942312)(322.4197263,286.28075072)(322.16953125,285.91161062)
\curveto(321.9275388,285.53836817)(321.8003901,284.99696192)(321.7921875,284.28739187)
\lineto(321.7921875,280.62879812)
\lineto(319.396875,280.62879812)
\lineto(319.396875,284.51707937)
\curveto(319.396875,285.34559552)(319.3230474,285.87879812)(319.18359375,286.10848562)
\curveto(319.0400388,286.34637572)(318.78984375,286.46942312)(318.42890625,286.46942312)
\curveto(317.96953125,286.46942312)(317.61269505,286.28075072)(317.3625,285.91161062)
\curveto(317.1082026,285.53836817)(316.98515625,285.00106322)(316.98515625,284.30379812)
\lineto(316.98515625,280.62879812)
\lineto(314.58984375,280.62879812)
\lineto(314.58984375,288.09364187)
\lineto(316.98515625,288.09364187)
\lineto(316.98515625,286.99442312)
\curveto(317.28046875,287.42098562)(317.6167974,287.73680567)(318.00234375,287.94598562)
\curveto(318.3837888,288.16336817)(318.7980474,288.27411062)(319.24921875,288.27411062)
\curveto(319.77421875,288.27411062)(320.23359375,288.14696192)(320.62734375,287.89676687)
\curveto(321.02109375,287.64247052)(321.32050755,287.29383692)(321.5296875,286.84676687)
\closepath
\moveto(321.5296875,286.84676687)
}
}
{
\newrgbcolor{curcolor}{0 0 0}
\pscustom[linestyle=none,fillstyle=solid,fillcolor=curcolor]
{
\newpath
\moveto(336.28710938,284.38582937)
\lineto(336.28710938,283.69676687)
\lineto(330.70898438,283.69676687)
\curveto(330.76230442,283.13895437)(330.96328177,282.71649317)(331.31601562,282.43348562)
\curveto(331.66464817,282.15868067)(332.15273437,282.02332937)(332.77617187,282.02332937)
\curveto(333.26835937,282.02332937)(333.77695312,282.09305567)(334.30195313,282.23661062)
\curveto(334.83515677,282.38836817)(335.38476562,282.61805567)(335.94257813,282.92567312)
\lineto(335.94257813,281.08817312)
\curveto(335.37246068,280.86668822)(334.80234427,280.70262625)(334.23632812,280.59598562)
\curveto(333.67851562,280.48934552)(333.11249947,280.43192312)(332.54648437,280.43192312)
\curveto(331.20117188,280.43192312)(330.15117187,280.77645437)(329.39648438,281.46551687)
\curveto(328.64999948,282.15457937)(328.28085937,283.11434552)(328.28085937,284.35301687)
\curveto(328.28085937,285.56707937)(328.64589818,286.52274317)(329.38007813,287.22411062)
\curveto(330.11015678,287.92137572)(331.12324193,288.27411062)(332.41523438,288.27411062)
\curveto(333.58417943,288.27411062)(334.51933567,287.91727442)(335.22070312,287.20770437)
\curveto(335.93027317,286.50633692)(336.28710938,285.56707937)(336.28710938,284.38582937)
\closepath
\moveto(333.84257812,285.17332937)
\curveto(333.84257812,285.63270437)(333.70312447,285.99774317)(333.43242187,286.27254812)
\curveto(333.16992188,286.54325072)(332.82539063,286.68270437)(332.39882812,286.68270437)
\curveto(331.93945312,286.68270437)(331.56210937,286.55145437)(331.26679688,286.28895437)
\curveto(330.97968697,286.03465802)(330.80742187,285.66551687)(330.74179687,285.17332937)
\closepath
\moveto(333.84257812,285.17332937)
}
}
{
\newrgbcolor{curcolor}{0 0 0}
\pscustom[linewidth=2.10000003,linecolor=curcolor]
{
\newpath
\moveto(77.540043,403.06453375)
\lineto(76.440903,315.87763675)
}
}
{
\newrgbcolor{curcolor}{0 0 0}
\pscustom[linestyle=none,fillstyle=solid,fillcolor=curcolor]
{
\newpath
\moveto(76.338276,308.00263675)
\lineto(71.223642,318.56826175)
\lineto(76.440819,315.87763675)
\lineto(81.723642,318.43701175)
\closepath
\moveto(76.338276,308.00263675)
}
}
{
\newrgbcolor{curcolor}{0 0 0}
\pscustom[linewidth=2.10000003,linecolor=curcolor]
{
\newpath
\moveto(76.338276,308.00263675)
\lineto(71.223642,318.56826175)
\lineto(76.440819,315.87763675)
\lineto(81.723642,318.43701175)
\closepath
\moveto(76.338276,308.00263675)
}
}
{
\newrgbcolor{curcolor}{0 0 0}
\pscustom[linewidth=2.10000003,linecolor=curcolor]
{
\newpath
\moveto(298.134375,401.54285275)
\lineto(297.863685,314.89326175)
}
}
{
\newrgbcolor{curcolor}{0 0 0}
\pscustom[linestyle=none,fillstyle=solid,fillcolor=curcolor]
{
\newpath
\moveto(297.839052,307.01826175)
\lineto(292.621875,317.53466275)
\lineto(297.863664,314.89326175)
\lineto(303.121875,317.50183975)
\closepath
\moveto(297.839052,307.01826175)
}
}
{
\newrgbcolor{curcolor}{0 0 0}
\pscustom[linewidth=2.10000003,linecolor=curcolor]
{
\newpath
\moveto(297.839052,307.01826175)
\lineto(292.621875,317.53466275)
\lineto(297.863664,314.89326175)
\lineto(303.121875,317.50183975)
\closepath
\moveto(297.839052,307.01826175)
}
}
{
\newrgbcolor{curcolor}{0 0 0}
\pscustom[linewidth=2.10000003,linecolor=curcolor]
{
\newpath
\moveto(518.724612,398.97117175)
\lineto(519.286572,315.33211975)
}
}
{
\newrgbcolor{curcolor}{0 0 0}
\pscustom[linestyle=none,fillstyle=solid,fillcolor=curcolor]
{
\newpath
\moveto(519.339849,307.45711975)
\lineto(514.020108,317.92020175)
\lineto(519.28653,315.33211975)
\lineto(524.520108,317.98992175)
\closepath
\moveto(519.339849,307.45711975)
}
}
{
\newrgbcolor{curcolor}{0 0 0}
\pscustom[linewidth=2.10000003,linecolor=curcolor]
{
\newpath
\moveto(519.339849,307.45711975)
\lineto(514.020108,317.92020175)
\lineto(519.28653,315.33211975)
\lineto(524.520108,317.98992175)
\closepath
\moveto(519.339849,307.45711975)
}
}
{
\newrgbcolor{curcolor}{0 0 0}
\pscustom[linewidth=2.10000003,linecolor=curcolor]
{
\newpath
\moveto(742.559181,398.76610675)
\lineto(740.959611,306.95673175)
}
}
{
\newrgbcolor{curcolor}{0 0 0}
\pscustom[linestyle=none,fillstyle=solid,fillcolor=curcolor]
{
\newpath
\moveto(740.959569,306.95673175)
\curveto(738.334569,307.00595575)(735.754677,309.67606375)(735.803901,312.30106375)
\curveto(735.848001,314.92196875)(738.519138,317.50183975)(741.144138,317.45673175)
\curveto(743.769138,317.41263175)(746.344914,314.74149475)(746.299806,312.11649475)
\curveto(746.255706,309.49149475)(743.584569,306.91160275)(740.959569,306.95673175)
}
}
{
\newrgbcolor{curcolor}{0 0 0}
\pscustom[linewidth=2.10000003,linecolor=curcolor]
{
\newpath
\moveto(740.959569,306.95673175)
\curveto(738.334569,307.00595575)(735.754677,309.67606375)(735.803901,312.30106375)
\curveto(735.848001,314.92196875)(738.519138,317.50183975)(741.144138,317.45673175)
\curveto(743.769138,317.41263175)(746.344914,314.74149475)(746.299806,312.11649475)
\curveto(746.255706,309.49149475)(743.584569,306.91160275)(740.959569,306.95673175)
}
}
{
\newrgbcolor{curcolor}{0 0 0}
\pscustom[linewidth=2.10000003,linecolor=curcolor]
{
\newpath
\moveto(76.309569,265.75653775)
\lineto(75.99375,179.40634375)
}
}
{
\newrgbcolor{curcolor}{0 0 0}
\pscustom[linestyle=none,fillstyle=solid,fillcolor=curcolor]
{
\newpath
\moveto(75.965043,171.53134375)
\lineto(70.751961,182.05183975)
\lineto(75.99375,179.40634375)
\lineto(81.251961,182.01082675)
\closepath
\moveto(75.965043,171.53134375)
}
}
{
\newrgbcolor{curcolor}{0 0 0}
\pscustom[linewidth=2.10000003,linecolor=curcolor]
{
\newpath
\moveto(75.965043,171.53134375)
\lineto(70.751961,182.05183975)
\lineto(75.99375,179.40634375)
\lineto(81.251961,182.01082675)
\closepath
\moveto(75.965043,171.53134375)
}
}
\end{pspicture}

%\caption{A simple block diagram\label{diagram}}
%\end{figure}

\end{comment}
